\documentclass[12pt,a4paper,openany]{ctexbook}
\usepackage{geometry}
\geometry{margin=1in}
\usepackage{hyperref}
\usepackage{setspace}
\onehalfspacing
\usepackage{fancyhdr}
\pagestyle{fancy}
\fancyhf{}
\fancyhead[LE,RO]{\thepage}
\fancyhead[RE]{\leftmark}
\fancyhead[LO]{\rightmark}
\renewcommand{\headrulewidth}{0.4pt}
\usepackage{indentfirst}
\setlength{\parindent}{2em}
\begin{document}

\title{帷幕之地 \cdot 灵焰纪元}
\author{帷幕之地世界组}
\date{\today}
\maketitle
\tableofcontents
\newpage

\begin{quote}
*\textbf{——奈莉娅，《帷幕观测录》第七卷}*

\end{quote}


\bigskip\hrule\bigskip


\chapter{卷一：太古纪元}


\section*{第一章：六座圣座}


\textbf{在那个人类尚以燧石取火、以结绳纪年、以洞穴深处的赭石手印确认"我存在过"的遥远纪元，诸神行走于帷幕之域的每一寸土地之上。}


\textbf{那不是虚无缥缈的信仰投射，不是原始心智对自然现象的拟人化解释。诸神真实存在。他们就住在凡人可以仰望、可以攀登、可以——如果你足够勇敢或足够愚蠢——叩门求见的宫殿里。霜脊山脉的玄冰巨城在正午阳光下的反光能让三百里外的旅人误以为看到了第二颗太阳；裂霆高地的银色圆顶在雷暴之夜像一颗从天穹落下的、被闪电环绕的星辰；幽翠圣林的圣橡巨树在灵质层面中的枝叶比任何石造宫殿都更为宏大壮丽，它们的根系不是扎在土壤里——是扎在一个凡人永远无法触及的、叫做"灵渊"的维度裂缝中。}


\textbf{但他们并非居住在同一个神域。在漫长的太古岁月中，六座圣座各自独立存在，各有疆域，各有律法，各有他们创世之初就承袭下来的、不可被外族理解的古老仪式。}


\textbf{在最北端，霜脊山脉的锯齿状峰顶刺穿了永冬的天空。那里屹立着}**\textbf{霜脊圣座}**\textbf{——一座由万古不化的玄冰与深黑色的玄武岩筑成的巨城。它的城墙高到可以将整片云层拦腰截断。它的塔尖越过大气中最冷的那一层，伸入星辰之间闪烁的虚空。在冬至之夜——当极光在霜脊上空展开它那年复一年不变但从未被任何人看厌的翡翠色和玫瑰色的光幕时——那些塔尖上的冰晶会被点亮，让整座圣城看起来像一座悬浮在天空中的、由冻结的星光塑成的宫殿。}


\textbf{霜脊诸神住在这里。他们的首领是独目者}**\textbf{瓦尔德林}**\textbf{。}


\textbf{在霜脊的语言中，"瓦尔德林"的意思不是"独目"——是"换取"。他年轻时是一位比其他任何神祇都更贪婪地追求智慧的神。当他听说霜脊山脉最深处有一口叫做"霜星之井"的泉眼——一口据说能让你理解万物运行的底层规律的泉眼——他没有犹豫。他独自走到那口被永冻土包围的井边，看了一眼井水中自己的倒影。那时的他还有两只眼睛。两只完整的、明亮如极昼的、能看穿九层迷雾的湛蓝眼睛。}


\textbf{他问井水："代价是什么。"}


\textbf{井水没有发出声音。但他在井水的倒影中看到了答案——他的右眼在倒影中自动闭上了。当他试图重新睁开它时，他发现他已经不能了。不是受伤，不是失明——是那只眼睛不再属于他。它被井水取走了。作为回报，井水将自创世以来封存在井底的所有智慧——总共是一万一千道"霜纹"——灌入了他的左眼。}


\textbf{瓦尔德林从此只用左眼看世界。而他的左眼——被那一万一千道淡青色的符文像星环一样环绕着——能看到过去九千年和未来九日。不能更远。九千年和九日，不多不少。他从未告诉任何人为什么是这两个数字。但在一个无人注意的冬夜，他在霜脊圣座的最高塔上站了整整九日。九日后他走下来，脸色比他上去时苍白了一千年。}


\textbf{他的右眼框从此空洞着。但他不遮它。他用左眼看未来，用空眼框记过去。}


\textbf{瓦尔德林有两个孩子。儿子叫}**\textbf{托尔迦}**\textbf{——雷霆锻铸者。托尔迦是六座圣座所有神祇中最强大的战士。他的风暴之锤"碎山者"——一块由一颗被冻结的恒星核心锻造的锤头，锤柄用的是瓦尔德林年轻时从霜星之井边折下的一根世界树祖枝——每一次挥击都能劈开山脉、轰出新的峡谷。风暴之锤在击中目标的瞬间会释放一道能照亮半个霜脊山脉的青白色闪电。闪电过后，空气中会残留三日的焦铁味和一种说不出的、像是冬雪和春雷同时存在的气味。托尔迦的战绩刻在霜脊圣座正殿的玄冰墙上——从左到右，总共是一千零二十三个被他击杀的太古巨兽和异域入侵者的名字。其中有三个名字被用红冰刻写（红冰是霜脊山脉中最稀有的冰种——只有在零下八百度以下才会形成的冰，它的红色来自一种被封闭在其中的远古藻类的灵焰残响）。问及为什么这三个名字是红的，托尔迦从不回答。}


\textbf{但他的父亲知道。那三个名字——在他杀死他们之后——再也没有任何人记得他们的存在。他们不仅是死了——他们是被从历史中抹去了。用红冰刻他们的名字，是托尔迦唯一能做的忏悔：他忘记了他们的脸，但记住了他们曾经存在。}


\textbf{瓦尔德林的女儿叫}**\textbf{芙蕾维娅}**\textbf{——华纳族的最后遗孤。华纳族曾是另一座比霜脊更古老、更柔和、更少涉入战争的圣座的守护者。那座圣座叫做"麦野圣地"——一个以金色麦浪、温暖海风和永不凋谢的野花著称的地方。那里的神祇不造武器、不修堡垒、不筑玄冰巨城。他们只种东西。种麦子，种果树，种记忆（是的，在麦野圣地，你可以把一段记忆种在土里，让它长成一株会散发那段记忆中的气息的花）。麦野圣地的存在维持了大约两万年。两万年后——在诸神之间的一场古老纷争中——它被从地图上彻底抹去了。不是占领，不是征服——是"抹去"。麦野圣地被一种无法被任何已知神术还原的力量从存在中删除。连带着它所有的人、所有的树、所有的花、所有的记忆。只有一个人活了下来——一个四岁的女孩，她当时正被她的父亲抱在怀里，她父亲在最后一刻将她推入了一道随机传送裂隙。裂隙的另一头是霜脊山脉的南麓，瓦尔德林在那里发现了她——蜷缩在一棵枯死的世界树祖枝下，嘴里含着一片金色的麦叶，在哭，但哭不出声。}


\textbf{瓦尔德林收养了她。她长大后成为霜脊圣座中唯一来自异族的神——她拥有华纳族的灵质感知力（一种能让灵焰触及其他生命灵焰的天然天赋）和霜脊的韧性（在永冬环境中的二十余年生活磨出来的）。她的灵焰是双色的——一半是霜脊的淡青色，一半是麦野的金黄色。没有人知道这意味着什么。芙蕾维娅本人也不确定。但她知道——每个夏天，当霜脊山脉短暂的融雪期到来，她会一个人走下山，走到山脚那片她能找到的最像麦田的草甸上，躺下来，闭上眼睛，试图回忆起那片金色麦浪的颜色、那片温暖海风的味道、那段两万年前在麦野圣地某个夏日黄昏中的、属于一个已经不存在的世界的宁静。}


\textbf{她一次也没有成功过。但她每年都试。她已经试了两千年。}


\textbf{霜脊圣座的守望者——}**\textbf{赫尔米昂}**\textbf{——站在虹桥的尽头。虹桥是一条连接霜脊与其他世界的光桥，由冰晶、灵质和一种叫做"晨光丝"的稀有材质编织而成。赫尔米昂的工作是守望——不分昼夜，不眠不休。他的金瞳能看到千里之外一只蜻蜓振翅的频率。他不是生来就拥有这双金瞳的。他年轻时曾是一个普通的哨卫——那时候他的眼睛还是普通的湛蓝色。在一次裂隙入侵中，他为了及时发出警报，在没有任何防护的情况下直接凝视了一道被帷幕之外力量撕开的裂隙。裂隙中翻涌的混沌物质在一瞬间灼伤了他的双眼——不是烧毁，而是将某种更古老的、从未被凡人感知过的光频段植入了他的视网膜。恢复之后，他的眼睛变成了金色。金瞳能看到所有肉眼不可见的灵质活动——每一个生命的灵焰的每一次涨落、每一条裂隙的每一个开合、空气中最微弱的灵质波纹。}


\textbf{赫尔米昂在这岗位上守望了三万天。他不睡觉——不是不需要，而是不敢。在他之前，有三位守桥人因为一瞬的疏忽而葬送了无数人的生命。他的名字"赫尔米昂"在霜脊古语中意为"永不阖目"。三万天——八十二年——他一次都没有阖过眼。他的金瞳始终睁开。但他会在每天黎明时分短暂地垂下眼帘——不是休息，是致哀。致那些他在三万天里看到的、从远方传来的、无法及时赶去救援的死亡。}


\textbf{向南走——越过霜脊山脉南麓那些藏着亿万年前就已灭绝的太古生物化石的峡谷，穿过那片被叫做"风嚎荒原"的、永远刮着能将岩石切成细沙的狂风的台地——你会进入一片由巨大裂谷和嶙峋岩峰构成的宏伟地域。这是}**\textbf{裂霆高地}**\textbf{。}


\textbf{裂霆圣座建在天穹之冠——最高那座峰——的峰顶。它的银色圆顶在雷暴之夜被闪电不断击打，但从不留下焦痕。裂霆圣座的主人叫}**\textbf{泽法利昂}**\textbf{，天霆尊主。他是六座圣座所有神祇中最威严的一个——也是最疲惫的一个。他的天穹之雷能击穿九层大地，他的雷矛——"天怒"——投出后能在云层中逗留三天三夜持续吸收电荷，第四天早晨带着足以轰碎一座中型山脉的能量自动回到他手中。}


\textbf{泽法利昂有两个兄弟和一个庞大的神系。他的长兄}**\textbf{哈德昂}**\textbf{统治着无光之渊——一个位于大地深处、从未有一丝阳光进入的黑暗王国。哈德昂不是邪恶之神——恰恰相反，圣座中少数明白"安静"之价值的神就是他。他所掌管的无光之渊是所有死者灵焰的最终收容所。死去的凡人和死去的神祇（是的，神也可以死去——只是比凡人慢得多，慢到有时候他们自己都不知道自己正在死亡）在无光之渊中继续着一种无需呼吸、无需心跳、但依然有意识的半存在。无光之渊不是惩罚，不是囚禁——它只是一间巨大的、安静的、黑暗的房间。那些在里面者的灵焰残响偶尔会哼唱从生者世界带来的歌。歌声在没有空气的黑暗中无法传播，但残响之间能以某种超越了声波的方式感知彼此的振动。他们称这种方式为"渊鸣"。}


\textbf{次兄}**\textbf{波里冬}**\textbf{统治沉没之殿——一片被无尽深水覆盖的古老宫室。波里冬的孤独是六座圣座所有神祇中最彻底的——不是因为没有别的神爱他，而是因为他无法离开深水。他出生时，一位早已被遗忘的远古神祇在他的灵焰中刻下了一道不可篡改的符文："永系于水"。离开水的波里冬会在大约一个时辰内开始灵焰逸散——不是受伤，是像一条被从鱼缸中捞出的鱼，在空气中逐渐窒息。所以在水面上，波里冬只有短短一个时辰。一个时辰，六十分钟，三千六百秒。他用来拜访他的兄弟。他每年出水一次——在冬至那一天，因为那是唯一一天泽法利昂会卸下天霆尊主的威严、变成一个愿意安静地坐在火堆边烤鱼的温和兄长的日子。兄弟三人围坐在裂霆圣座正殿的铜炉边。哈德昂无法来到地面上——他不能离开无光之渊超过一盏茶的时间，因为无光之渊的"无光"环境是维系他灵焰稳定性的必要前提。所以他和波里冬之间的对话是通过泽法利昂来转达的——泽法利昂在两个兄弟之间充当了一万多年的传话筒。}


\textbf{南方的幽翠圣林是翠庭圣座的所在。这片森林中的每一棵树都是活的——不是普通意义上"有生命的"，而是"有意识的"。圣橡的根系在灵质层面中交织成一张覆盖整片林地、延伸到无法测量距离的神经网络。这张网络叫做"根语"，翠庭诸神通过它将思维、记忆和梦境在彼此之间共享。根语让翠庭拥有了六座圣座中最快的通信速度——一个神在林地北端看到的东西，南端的所有神在同一瞬间也看到了。速度——比闪电、比光、比任何已知的信息传播方式都快——是即时的。}


\textbf{翠庭的首领是}**\textbf{达格德里安}**\textbf{，圣父。他是所有圣座中最不像"统治者"的神——他从来不坐在王座上，从来不发布敕令，从来不要求任何人称他为"尊主"或"圣父"。他每天做的事情是：在黎明前起身（他不需要睡眠，但他喜欢在清晨的寂静中思考），走遍幽翠圣林的每一株圣橡，将手按在每一株老树的树皮上，问它们"昨夜风大吗？"。他会得到所有的回答。不是用语言——根语传递的信息比语言更丰富也更高效。他会在掌心感受到每一株树在过去一夜中经历的每一次风的拉扯、每一滴露水的滑落、每一只夜行动物在枝头轻踏时的振动。三息之内，他知道整片森林昨夜发生的一切。这是他的一日之始。}


\textbf{翠庭诸神中有一个女神的名字在六座圣座中无人不晓——}**\textbf{莫丽甘娜}**\textbf{，命运与渡鸦的女神。她能以渡鸦的眼睛看世界——不是一只渡鸦，而是翠庭领空内每一只渡鸦同时成为她的眼睛。所以她同时看到翠庭的每一个角落——南边溪流中一条跃出水面的鳟鱼、北边一棵圣橡上新生的嫩芽、西边草丛中两只争夺领地的刺猬、东边一棵枯树上最后的黄叶离开枝头的那个瞬间。}


\textbf{最令人敬畏的是——她记住了。她记住了自她诞生以来通过所有渡鸦看到的所有瞬间。她的大脑——不，她的灵焰——中储存了翠庭一万八千年来的完整视觉档案。一万八千年——比大多数文明的存在时间都长。她可以随时调取其中任何一帧画面——比如第四千三百二十二年春分日黄昏、第十七秒、南偏西那只渡鸦看到的、一只雌鹿低头饮水的侧影。她能原样重构。}


\textbf{但她最令人不安的还不是这个。她偶尔会以渡鸦的形态出现在战场上——不是参与战斗，而是等待战斗结束。当战场上最后一声嘶喊消逝、最后一面旗帜倒下——莫丽甘娜的渡鸦会落在一个垂死战士的盔甲上。她会以渡鸦的形态对他说一段话。没有人知道她说的内容——每一个听到这段话的人都在片刻后死去了。只知道那些死者的脸上——每一个人的脸上——都带着一种无法描述的表情。不是恐惧，不是痛苦。是满足。像是在最深的黑暗中，有人终于、终于、终于解开了他们一生中最困惑的谜题。}


\textbf{最南端的圣座在金沙之畔。那里曾经是——现在仍然是——一片被金色沙粒覆盖的广袤平原。沙粒不是岩石风化形成的。它们是亿万个被遗忘的名字。}


**\textbf{金沙圣座}**\textbf{——日舟之主}**\textbf{拉赫卡里昂}**\textbf{的领地。拉赫卡里昂每一天驾驶他的太阳舟从东天航行到西天，将光和热带给帷幕之域的每一个角落。这是一个看似辉煌、实则令人筋疲力尽的仪式。太阳舟不能停下——不是因为有谁强迫他，而是因为太阳舟本身就是一个活的、自主的存在，而它与拉赫卡里昂之间的契约规定：只要拉赫卡里昂还有一口气，太阳就必须升起。这个契约是他出生时由他不知道的远古存在刻下来的。他不能解约。他不知道解约的条件。他甚至不知道该向谁提出解约。所以他每天划动太阳舟。他已经划了近两万年。他的手从来没有停过。唯一让他继续下去的动力，是他知道——如果他停下来，那些在黑暗中恐惧的孩子就永远看不到黎明了。}


\textbf{拉赫卡里昂的伴侣是}**\textbf{伊希丝}**\textbf{——万符之母。伊希丝是六座圣座所有神祇中最精通灵焰文字的神。她所创造的"灵符"是一种能将灵焰信息编码为符号的系统——每一个灵符代表着一个特定的灵焰状态、一个特定的情感色调、或一个特定的记忆片段。她将这些灵符刻在骨灰平原上的白石板上，用来记录金沙圣座第一万零一年起的完整历史。至今已有九千多块石板排列在金沙圣座的地下档案馆中——一家叫"命符殿"的、位于地下三里的、温度始终恒定为某个特定度数的、除了伊希丝本人没有人进去过的巨大石室。}


\textbf{她的儿子}**\textbf{安努赫卡}**\textbf{——冥途秤官——是所有死去灵魂的审判者。每一个亡灵在进入无光之渊之前，都必须经过安努赫卡的天平审判：将亡者的灵焰置于天平的左侧，将一根来自圣橡的纯白羽毛置于右侧。如果灵焰因善行而轻盈、因罪恶而沉重，天平就会倾斜。根据倾斜的方向和幅度，安努赫卡决定这个亡灵的去向：沉入无光之渊的深处（漫长而静止的半存在）或升入日舟的光辉中（短促但炽烈的二次燃烧）。}


\textbf{安努赫卡的审判从未出错。但他在每次审判结束后——在每一声"向左"或"向右"的宣判之后——都会独自在他的天秤宫里坐一炷香的时间，一遍又一遍地重新过一遍审判过程。他怕错判。他说服不了自己他的审判永远正确。他在自己给自己设置的上诉期内一遍又一遍地重新思考每一个亡灵应该去的地方。他重审了整整一万六千年。至今没有推翻过一次原判。但他仍然每天重审，一天不落。}


\textbf{西北方向——越过那片被称为"千界回廊"的、由被遗弃的古代传送门碎片拼接而成的荒原——有一座用纯白大理石建造的巨城。这就是}**\textbf{白石圣座}**\textbf{——律法与秩序的圣城。白石圣座的主人叫}**\textbf{朱沃纳斯}**\textbf{，律法之父。朱沃纳斯最大的贡献——在他的所有成就中——是编纂了《万域通则》：一套由一千一百位法官和学者在两百年间起草、由朱沃纳斯亲自修改了七十八稿的法律典章。万域通则规范了六座圣座之间的所有潜在纠纷：边界划定、贸易、跨界婚姻（当一位霜脊神祇爱上一位翠庭神祇时，他们该住在哪座圣城？——答案是他们可以在两座圣城之间选一座，但不能在两个地方同时拥有神殿。这条规定至今被引用过四十七次）。万域通则是六座圣座之间之所以能维持万年和平的制度基石。但朱沃纳斯从未告诉任何人——在通则的最后一页、最后一行的最后一个字后面——他亲手用只有在特定月色下才能看到的墨水加了一个字。那个字不能被直接阅读。它的存在——它在那里的这个事实本身——才是最重要的：最后一个字是一个没有被写出的字，一个空白的印记，一个被故意的、留白的、也许永远不需要被填上的"如果"。}


\textbf{白石圣座最神秘的神祇是}**\textbf{雅努维翁}**\textbf{——双面门神。他的两颗头分别朝向过去和未来。他建造的门不是物理的门——是时间的门。在帷幕之域的各个角落，他留下了不为人知的、通往特定时间点的小门。没有人知道这些门的位置——连雅努维翁自己也不记得了（因为他的两颗头每两千八百年会自动重置一次记忆，旧记忆被送入无光之渊储存，新记忆从空白开始）。如果有人偶然进入了其中一扇门，他们会出现在门通向的时间点——可能是三千年前的某个春日，可能是三千年后的某个冬夜。门的规律无人能解。唯一知道的是：所有通过雅努维翁之门的人，都没有回到他们原来的时间。}


\textbf{然后是东方的}**\textbf{雾霭圣座}**\textbf{——一座悬浮在永恒薄雾中的仙岛群。每一座仙岛由一道只有特定步幅才能走的石径连接——错一步就踏空，跌入永远也落不到底的白雾中。雾霭圣座的尊主是}**\textbf{玉穹尊者}**\textbf{——一个从不展示真容、只以声音和雾中显形来交流的神祇。他是六座圣座中最难以理解的一位——不是因为他表达不清楚，而是因为他表达的方式不是语言。他通过将思维以"意境"的形式直接投射到对方灵焰中。你无法复述玉穹尊者说了什么——你只能说你感受到了什么。而每个人感受到的都不同。所以雾霭圣座从来没有产生过统一的教义——不需要——因为每个人与玉穹尊者之间的交流都是私密的、不可复制的、只存在于那一瞬灵焰共振之中的。}


\textbf{雾霭圣座还有一位叫做}**\textbf{天照姬}**\textbf{的女神——她是六座圣座所有女神中最沉静的一位。天照姬掌管旭日的第一道光。她在日舟升起之前，先将自己灵焰中的一缕热和一分亮投射到东方的天幕上——让凡人在太阳冒出地平线之前先感受到"新的一天将要开始了"的信号。她从不说话。但她每天早晨的这一个小动作——这个持续了大约十六秒的灵焰投射——是所有凡人一天中最不害怕的那十六秒。他们说，在那一刻，即使你昨晚失去了最重要的东西——那十六秒的微光依然会让你觉得：也许今天还可以再试一次。}


\textbf{六座圣座共存了数万年。数万年里，他们之间有过战争，有过联姻，有过长达几千年的冷战，也有过让所有凡人都暂时松了一口气的千年和约。但不管他们之间发生过什么——他们都共享一个边界。}


\textbf{那个边界没有名字。给它取名本身就会赋予它不应拥有的力量。}


\textbf{六座圣座的智者们各自用自己的方式称呼它。霜脊称之为"深空之外的深空"——那个比星辰之间更远的地方。裂霆称之为"雷霆的尽头"——泽法利昂的雷矛投到最远也无法触及的区域。翠庭称之为"根语中的沉默"——圣橡的神经网络无法感知到的、空白地带的上方。金沙称之为"日舟不渡之水"——拉赫卡里昂的太阳舟永远不经过的那片虚空。白石称之为"通则之外的条款"——朱沃纳斯亲口承认他的法律无法管辖的领域。雾霭称之为"雾中的不是雾"——当你穿透最浓的雾之后，你看到的那个东西，不是另一层雾。}


\textbf{而所有的称呼——不管有多么不同——都指向同一个真相：}


\textbf{在那里，存在着比诸神更古老、更饥饿、更不会死去的东西。}


\section*{第二章：第一道裂隙}


\textbf{裂隙最初出现在霜脊山脉北麓一个无人在意的冰洞里。}


\textbf{那个冰洞连名字都没有。它太小了——洞口窄到只容一个人侧身挤入。洞内的空间大约只有一间小储藏室大小。唯一的特征是洞壁上的冰是深蓝色的——比其他地方的冰蓝得多，蓝到近乎发黑。老人们说洞里有"冰下的光"——在特定的凌晨，当月光以一个精确的、每年只有三天的角度射入洞口时，洞壁上的深蓝冰会发出一种微弱的、不断呼吸的磷光。}


\textbf{在裂隙出现之前，那只是一个普通的奇异冰洞——奇异，但无害。}


\textbf{但在那一年——霜脊圣座编年史上称之为"无光之年"——赫尔米昂的守望记录中出现了一行他从未见过的东西。}


\textbf{那天是冬至前第七日。赫尔米昂照例站在虹桥尽头，用那双永不阖目的金瞳扫视九层天穹。他看过了霜脊山脉——平安。他看过了裂霆高地——平安。他看过了幽翠圣林——南侧的圣橡在昨夜的风中落了一根百年枯枝，但这不是威胁。他看过了金沙——日舟正在按照精确的航路行驶，拉赫卡里昂的划桨频率和昨天一样，比两万年前略慢了一拍（他老了，赫尔米昂心想。我们都老了）。}


\textbf{然后，出于某种他后来无法解释的冲动——也许是金瞳在长年守望中形成的一种直觉——他看向了一个他从未看过的地方。}


\textbf{那是极北之北。比永冬山脉更北。比风的起源地更北。比光能够抵达的最远点还要北。在那个方向，在物理空间的意义上，不应该存在任何东西。帷幕之域在此处达到了它的最大外延——像一张铺开的羊皮纸在此处结束了它被墨水覆盖的最后一丝痕迹。外面——如果在"外面"还有东西的话——是纯粹的混沌：没有光，没有物质，没有时间，没有空间——没有"有"。}


\textbf{赫尔米昂的金瞳看到了那里。}


\textbf{然后他看到了那个东西。}


\textbf{他僵在虹桥尽头——不是恐惧（他是守桥人，三万天来他已经见过九千种威胁，他不会再恐惧了），而是无法理解。他在看到那个东西的同时意识到了一件事：他的金瞳在被灼伤时获得了看见灵质的能力。而灵质——是这个世界上一切存在物的本质。他能看见灵质，所以他能看见一切存在物。但他现在看到的东西——他意识到了第二件事——没有灵质。不是灵质是黑色的（那只是蚀痕），不是灵质是微弱的（那只是一息尚存者）。是——没有灵质。而一个没有灵质的东西，是不应该存在的。万物皆有灵焰。没有灵焰的东西就是不存在的——这是世界的基本公理。}


\textbf{他看着那个不存在的东西向帷幕之域蔓延过来。它的速度非常慢——慢到如果不是赫尔米昂一直在看着它，他根本不会察觉。它像一滴墨水落在湿润的羊皮纸上——从落点开始，向周边缓缓渗透。不——不是墨水。是"没有墨"。它是墨水的反面——空白。它不是在写字，而是在抹去已经被写下的字。它不是在添加——它在删除。}


\textbf{赫尔米昂拉响了虹桥上的警报钟。那座钟在霜脊圣座的中央大殿中响了十二声——这是最高等级警报，自从它被铸成以来只被拉响过三次。第一次——四万年前——是一头从星辰间游来的虚空巨兽。第二次——一万八千年前——是那次导致了麦野圣地覆灭的古老纷争。第三次——就是现在。}


\textbf{霜脊诸神在十八息之内全部赶到了。瓦尔德林用他那只能看到过去九千年和未来九日的左眼看向极北之北。他看了很久——比平常久得多。他的左眼中——环绕着一万一千道霜纹——每一道霜纹都在微微发光。}


\textbf{然后他阖上左眼。睁开空着的右眼框。用那个什么也看不见的眼框——"看"向同一个方向。}


\textbf{他的右眼框是一种工具。它看不见光，看不见灵质，看不见未来，看不见过去。但它能——瓦尔德林从未向任何人解释过这一点——感知到不存在。他的空眼框在感知不存在这方面的敏感度，比他那只被一万一千道霜纹加持的左眼精准一百倍。}


\textbf{他用空眼框对着极北之北——对着那个赫尔米昂看到的方向。一动不动。}


\textbf{过了很久。太阳舟从东天划到了西天，又从东天重新升起了一轮。霜脊诸神在虹桥尽头站了一整天。}


\textbf{最后，瓦尔德林空了的那只眼框——在所有人都没有预料到的时候——流下了一滴泪。不是悲伤的泪——是疼痛的。像一个空杯子被突然倒入了太烫的水。}


\textbf{"不是不存在，"他说。"是尚未存在。"}


\textbf{他转向赫尔米昂。}


\textbf{"你看到的不是'无'。你看到的——是'还没有'。那个东西还没有成为任何可以被称为东西的存在。但它正在学习。它以我们无法想象的缓慢速度渗透进来——它在观察。在品尝。在理解什么是'存在'。当它理解的那一刻——当它学会就说'我存在'的那一刻——它就不再是'还没有'了。它就是'是了'。"}


\textbf{赫尔米昂没有说话。但他的金瞳一直在看着那个逐渐逼近的——尚未存在的东西。他想问一个问题，但他知道他最不想听到瓦尔德的答案。}


\textbf{那个问题是：当它学会说"我存在"之后——它会做什么？}


\textbf{他最终没有问。因为他的金瞳已经给出了答案。在极北之北，在不应该有任何东西的地方——他看到了那个尚未存在之物的——尚未存在的——嘴。}


\textbf{正在张开。}


\section*{第三章：不可能之战}


\textbf{战争的第一枪不是神开的。}


\textbf{是那个尚未存在的东西——在它终于学会了"存在"之后——开的第一枪。}


\textbf{它现在已经不能被称为"尚未存在的"了。它已经学会了。它已经成为了。在瓦尔德林拉响警报之后的第四十九天，赫尔米昂的金瞳看到一个形态——如果他必须用语言描述的话——从极北之北的虚空中凝聚出来。那不是任何已知生物的形态。它没有固定的轮廓——它像一团不断翻滚、不断变形、不断在"某个形态"和"不准确的另一个形态"之间切换的黑雾。你盯它盯得越久，你的眼睛就越疼——不是生理上的，而是一种认知层面的过载：你的大脑试图理解一个不可理解的东西，而理解的过程本身就是痛苦的。}


\textbf{它会模仿。当你看着它时，它会变成你潜意识深处最恐惧的那个东西的形状。赫尔米昂看到了一个正在溺毙的孩子（他小时候在冰湖上滑冰时掉入冰窟——他至今仍会在梦中反复经历那件事）。泽法利昂说它看起来像一道从天而降、但永远不击中地面的雷——被悬在空中，无处释能。拉赫卡里昂说它看起来像一片黑暗的、没有日出的东天。而莫丽甘娜——当她通过一只渡鸦的眼睛看到它时——她说了所有人都没有预料到的话："它看起来像我。"}


\textbf{瓦尔德林决定出击。他不能让这个刚学会了存在的东西继续留在帷幕之域的边界上。他召集了霜脊的战士——托尔迦带着风暴之锤，赫尔米昂带着他的金瞳，芙蕾维娅带着她一半霜色一半金色的灵焰。}


\textbf{在霜脊山脉北麓——那道最薄的冰裂缝处——诸神第一次正面面对了那个东西。}


\textbf{托尔迦没有给它时间变形。他挥动风暴之锤，一道青白色的雷光击中那个东西——击散了它。就像水被劈开——黑色的雾气在雷光中向四周飞散。但飞散没有持续。在雷光消退之后，飞散的黑色雾气开始缓慢地——仿佛在思考、在计算、在分析——重新合拢。三息之内，它重新聚成了与之前一模一样的那团翻滚的黑色形态。}


\textbf{托尔迦又挥了一锤。散开。重新合拢。他连挥了十一锤。十一次——每一次都精确地打散，每一次它都重新合拢。第十一次合拢的速度比第一次快了一点。第十二次——锤还来不及挥——那个东西已经合拢了。它学得很快。太快了。}


\textbf{泽法利昂释放了他的天穹之雷——裂霆圣座的最强攻击。闪电击穿了黑雾——这次不止是打散，而是把黑雾的中心烧出了一个洞。一个完美的、圆形的洞。但洞在自我修复——像是伤口边缘的黑色物质在重新生长。泽法利昂持续放电，雷光在洞中继续燃烧——三息，五息，十息——洞开始缩小了，黑雾的边缘开始向中心涌——然后——在一瞬间——洞完全闭合了。雷光被吞没在了黑色的物质中。}


\textbf{泽法利昂的天穹之雷——可以把一座山脉轰碎的攻击——被吃了。}


\textbf{诸神沉默。}


\textbf{然后瓦尔德林做了一件所有人都没有预料到的事。他向前走了一步——走进了黑雾中。他走到了那个连他都不确定是什么的存在之中。他用他那只被一万一千道霜纹环绕的左眼近距离地看着它。他的右眼框——空的那只——在黑色物质的包围中突然剧烈地疼痛，他蹲下来，把空的眼框埋在双手之中。过了很久。他站起来，走回诸神身边。}


\textbf{"你们都猜过为什么我看不到九千年之后的未来，"瓦尔德林说。"我看了。那里什么都没有。不是'没有灾难'——是'什么都没有'。"}


\textbf{他停顿了一下。}


\textbf{"九千年之后这个世界就不存在了。"他继续说道。他的声音很平静——像在陈述一个已经发生的事实。"除非。"}


\textbf{"除非什么？"托尔迦问。}


\textbf{"除非我们把它挡在某个东西外面。一种——屏障。不是墙。墙会被它吃掉。是一种——"瓦尔德林停顿了一下，在霜脊的古语中搜寻一个合适的词。"——一种帷幕。不是挡住它，而是让它觉得不值得翻过来。让它觉得这边没有它想要的东西。让它——对它认知中的'存在'产生怀疑。"}


\textbf{他环顾诸神。}


\textbf{"你们需要帮我做一件事。不是你们想做的事。不是你们训练过做的事。是你们从来——从来没有想过自己会做的事。"}


\textbf{他没有说"那是什么"。但他左眼中的一万一千道霜纹在此刻全部熄灭了一瞬间——像暴风雨中被闪电短暂照亮的黑暗水面——然后又重新亮起。霜纹重新亮起的时候，每一道霜纹的中心都被染了一层很浅很浅的——如果不仔细看几乎看不出的——灰。}


\textbf{芙蕾维娅注意到了那个灰色。她后来在日记中写道："他在那个黑雾里面看到了什么——他不肯告诉我们。但那个东西把某种颜色留在了他的霜纹里。也许是悲哀。也许是计划。也许——是一个交易。我不确定。我从来都不确定。"}


\section*{第四章：会盟与分歧}


\textbf{瓦尔德林向其他五座圣座派出了信使。}


\textbf{这不是一件容易的事。六座圣座之间虽然有《万域通则》维持着和平，但那是脆弱的和平——是通过回避某些"不方便讨论"的话题来维持的和平。上一次六座圣座全体会晤还是在六千年前——那次的议题是"是否应该允许神与凡人通婚"（结果是维持现状：不鼓励，不禁止，不承认后代的圣座继承权）。这一次的议题完全不同。这一次的议题是：是否所有圣座愿意牺牲自己来编织一道帷幕。}


\textbf{信使们花了三十多天在各个圣座之间奔波。}


\textbf{裂霆圣座最先回应。泽法利昂在瓦尔德林的信使到达裂霆之前就已经准备好了——他的天穹之雷在数日前感知到了那个东西。他的天穹之雷不仅是一把武器——它还是一个巨大的灵质探测系统。它能感知到帷幕之域边缘的每一次灵质波动。最近几日的波动——泽法利昂在看了读数之后沉默了至少半个时辰——已经超出了他探测系统的量程上限。他的探测系统量程上限是"一千个太古巨兽"。读数——是"无法计量"。}


\textbf{哈德昂——泽法利昂的长兄、无光之渊的主宰——也做出了他自己的评估。他没有探测系统，但他有一种更直接的方法：他查看最近一个月内从生者世界进入无光之渊的亡灵。正常的死者——那些寿终正寝、灵焰自然熄灭的人——进入无光之渊时是安静的、有序的、像一滴水融入一口深井。但最近一个月——哈德昂的渊鸣中听到了不寻常的、不属于正常死者的"噪音"。那不是安静的融入——是暴力的撕裂。有什么东西——在他的渊之外——在强行把灵焰从活人体内剥离。那些被剥离的灵焰没有进入无光之渊。它们消失了。不是"去了别的地方"——是消失了。在灵焰层面上彻底消失。}


\textbf{哈德昂把这份报告用金箔写下，让一只可以在水下和空中都存活的古老生物——一种叫做"银鳐"的、他养了七千年的宠物——送到了裂霆圣座。报告末尾只有一行字：}


\textbf{"它在吃灵焰。不是像我们——我们安顿它们。它不是。它在吃。被它吃掉的灵焰——不再存在。没有痕迹。无光之渊里找不到它们。"}


\textbf{翠庭圣座第三个回应。达格德里安的态度可以用一句话概括：幽翠圣林的一切——每一棵圣橡、每一片苔藓、每一只在林地中鸣叫的蝉——都在对他说同一件事。"它们在恐惧。它们一直在恐惧。它们以为我没有注意到。但我注意到了——它们的根语中出现了同一个信号。反复地。一遍又一遍地。那个信号不是警告——是告别。它们在告别。它们以为我们不会保护它们。"}


\textbf{莫丽甘娜用她的渡鸦之眼巡逻了帷幕之域的整片北边界。她看到了和赫尔米昂相同的东西——但它更大了。比赫尔米昂最初报告时大了至少四倍。扩散速度也在加快。按照她的估算——如果扩散速度不变——在六年之内，黑雾将覆盖整个霜脊山脉。在十二年之内——它将达到幽翠圣林。}


\textbf{六年和十二年。在神的尺度上——那是弹指一挥间。}


\textbf{金沙圣座的回应稍慢——因为伊希丝坚持要在提交任何决定之前先完成一项实验。她用她发明的灵符——那些能将灵焰信息编码为符号的装置——来测试黑雾样本。拉赫卡里昂冒险在日舟靠近极北之北的航段中，用一种叫"永恒之冰"的容器（霜脊圣座之前送给金沙的礼物）捕获了极少量的——大约一撮盐的量——黑雾。}


\textbf{伊希丝对微量的黑雾做了她能想到的所有测试。她用了灵符，用了天平，用了日舟上的圣火，用了骨灰平原上所有已知物质的样本。所有测试都指向同一个结论——所有反应都呈现出同一个特征：黑雾不与任何东西发生反应。不是惰性——是拒绝。黑雾拒绝与你互动。拒绝与你反应。拒绝被你测量、被你知道、被你所拥有任何关于它的知识。如果你持续尝试——它会开始发出一种极低频的震动。那种震动——伊希丝在实验记录中写道——"让我想睡觉。不是困了。是那种——当你面对某个你永远无法对解答的问题时——你会觉得睡眠是唯一的安慰。因为只有在梦境里，你才不需要问那个问题。"}


\textbf{她的儿子安努赫卡——冥途秤官——做了最后一项测试。他将黑雾样本放在他审判亡灵所用的天秤上。用天平来测量灵焰的重量是他的职业本能。但黑雾——一个没有灵焰的东西——在天平上会有多重？答案让安努赫卡的面色在那一瞬间变得比骨灰平原的沙土更灰白。}


\textbf{"没有重量，"他说。}


\textbf{"没有重量——是零？"}


\textbf{"不是零。"安努赫卡停顿了很久。"是负数。不是质量是负的。是在它的存在之上——它还需要更多存在来抵消它的负存在。它在吃灵焰——因为它需要正面的'存在'去填补那个负的。它是——一个存在的缺口。一个永远填不满的空。一个饥饿的零。"}


\textbf{白石圣座接到信后的反应最为独特。朱沃纳斯用了一个多礼拜的时间研读《万域通则》——试图在他的法律中找到任何关于"面对一个不可战胜的存在、诸神应该怎么办"的条款。没有任何条款提供这方面的建议。万域通则是和平时期的产物——它假定宇宙是稳定的、规律是可循的、对手是可以在某种意义上与之谈判的。它从来没有考虑过——朱沃纳斯在读完最后一页后合上厚厚的法典，对他的书记官说——"一个连存在都还不确定的东西。"}


\textbf{他做了一个决定：在白石圣座的正殿中召集所有还在世的法官，公开讨论是否应该参与瓦尔德林的计划。讨论持续了六天六夜。正反两方僵持不下。最后——在第七天的凌晨——朱沃纳斯请出了雅努维翁。}


\textbf{雅努维翁的双面门神身份使他成为唯一一个能同时对过去和未来"说话"的神。不是预测——是实际地与过去和未来的自己进行某种程度的交流。雅努维翁用他面对过去的那颗头向后看——看到了过去一万八千年中白石圣座曾经面对过的所有危机。然后他用面对未来的那颗头向前看——看到了他没有说话。}


\textbf{他没有说话。}


\textbf{不是不想说。是不能说。他的未来头颅所看到的东西——在他开口准备描述的一瞬间——从内部被抹除了。像有人在他脑子里放了一块橡皮擦。他再试，再擦。第三次，未来头颅的眼眶中流出了血——不是因为受了伤，而是因为某个东西不允许那个未来被任何人看到。}


\textbf{朱沃纳斯扶住了即将倒地的雅努维翁。他在雅努维翁的耳边轻声问了一个问题。没有人知道那个问题是什么。雅努维翁用极其细微的声音回答了他。}


\textbf{朱沃纳斯听完后，站起来对所有在场的法官说了六个字：}


\textbf{"我支持。没有补充。"}


\textbf{雾霭圣座的玉穹尊者——和往常一样——没有做出可以被翻译为语言的任何回应。他只向瓦尔德林投射了一个意境。瓦尔德林在接收到这个意境的瞬间——他那只剩下一只的左眼中——落了一滴泪。不是疼痛的泪。是感激。没有人知道玉穹尊者到底"说"了什么。瓦尔德林也没有透露。但在他接下来的所有行动中——直到编织帷幕的那一天——他的左眼中多了一道之前从未有过的霜纹。那是一道青色的、但边缘闪着柔和的乳白色微光的霜纹。每当他在黑暗中闭上左眼，那一道霜纹会独自亮一小会儿，像一个遥远的、仍在等待归人的灯塔。}


\section*{第五章：诸神的集体抉择}


\textbf{最终的会盟地点选在了裂霆高地的天穹之冠——泽法利昂的主殿。那座银色圆顶下的大殿足以容纳上千位神祇。那天到场的没有上千——只有不到百位。幸存者们稀稀落落地站在宏伟得有些嘲讽的大殿中，彼此之间留着一种不自然的间距——不是疏远，是疲惫。太疲惫了，已经无法像万年前那样挤在一起、用彼此的体温来互相安慰。}


\textbf{瓦尔德林站在大殿中央。他的肩上停着最后一只渡鸦——不是莫丽甘娜的化身，而是他的老友穆宁（在霜脊的古老语言中意为"记忆"）的最后一羽。他的另一只渡鸦——胡金（"思想"）——在第三次裂隙侦察中飞得太近了，被黑雾吞没了。胡金被吞之前的最后一幕——通过它与瓦尔德林的灵焰链接传到瓦尔德林脑中的画面——是一片漆黑。不是看不见任何东西的黑——是看得见的黑。是一片有意识、有目的、正在看着胡金而胡金也正在被它看的黑。}


\textbf{瓦尔德林从那以后再也没有召唤过新的渡鸦。他说他不确定被黑雾吞没的胡金还在不在——如果它还在，他不想让另一只渡鸦也经历同样的事情。如果它已经不在——他不确定自己想不想知道。}


\textbf{"我不想用长篇大论浪费你们的时间，"瓦尔德林对殿中诸神说。"你们都已经看到了。你们也都已经知道——凭我们的武力，我们打不赢。它不能被杀死。不是因为它太强——是因为它没有可以被杀死的东西。它没有灵焰。没有心脏。没有可以被定义为'致命点'的解剖结构。它不是活的——它只是'在了'。而我们目前已经发现的事实是——你打它，它会学会你的打法。你用雷，它会吃雷。你用火，它会吃火。你越是攻击它，它越是理解你。它越理解你，它就越是——存在于你之中。它的存在是以我们的攻击为养料的。"}


\textbf{他停顿。殿中只有此起彼落的呼吸声（那些还可以呼吸的神在呼吸）和沉默（那些不需要呼吸的神在以沉默互相交谈）。}


\textbf{"所以我们能做的是另一件事。"}


\textbf{他抬起右手。他在掌心凝聚了一小团淡青色的光——那是他灵焰中抽出的一个碎片。整个大殿中温度骤降。不是物理上的——是灵质上的。每一个神的灵焰在被瓦尔德林的灵焰碎片照射时，都轻轻震了一下。像是被问到同一个问题：你准备好了吗？}


\textbf{"我们可以把自己变成一道屏障。不是挡住它——是封住它。把我们所有残存的神力——我们所有的灵焰碎片、我们所有的记忆、我们所有的爱、恨、悔、誓——编织成一道帷幕。一层一层的。最外层——由我们中的最强者在完全清醒的状态下将自己献祭为幕料。中间的层——由我们的记忆、我们活过的痕迹、我们爱过的面孔——那些无形的但更坚固的——作为幕心。最内层——由我们所有的希望——对——希望——那种最不被当作武器但却是唯一能在完全黑暗的环境中维持自身不灭的元素——作为幕底。"}


\textbf{他放下了手。}


\textbf{"这会让我们不复存在。不是死——你们听到了。死亡是将灵焰送入无光之渊——哈德昂会接待你。这个不是死亡。这是'不再作为一个独立意识存在'。你的灵焰会被打散、拉伸、编织成一层只有半个原子厚的薄膜。你不会再有'我'了。你不会再有思想、记忆、感受——你会变成一道屏障的一部分——没有意志，没有欲望，没有——'你'。"}


\textbf{长久的沉默。然后——令所有人意外地——第一个走向瓦尔德林的不是霜脊的任何一位神。而是伊希丝。}


\textbf{伊希丝从金沙圣座的诸神中走出来，手臂上刻满了一万道微小的金色符文。那是金沙圣座两万年来的全部历史——她用了三十天的时间将它们压缩到她的皮肤上。那些符文的体积不足以保存每一条记录的完整文字，但它们保留了每一条记录的灵焰签名：记录中涉及的人、事、年代、来龙去脉——所有关键信息都以一种只有伊希丝能解读的编码方式封存。}


\textbf{"我准备好了，"她说。"不是因为我最勇敢——是因为我最清楚我们欠那些还没有出生的人什么。安努赫卡告诉过我——在无光之渊里，他听到了婴儿的声音。那些婴儿还没有出生。他们还是灵焰——在他们母亲子宫里的那一小团还没有成形的灵焰。这些婴儿在无光之渊中发出了同一个问题——不是用语言——我不能翻译。但安努赫卡说那是一个问题。而那个问题——如果我们不编织这道帷幕——将永远没有人能回答。"}


\textbf{她走向大殿中央。}


\textbf{"我做这件事是为了那些还没有出生的人——那些还没有被问问题的人。让他们有一天能够问出他们想问的——能够等到有人给他们答案。"}


\textbf{然后是托尔迦。他放下风暴之锤。那把陪伴了他一万五千年的锤子躺在天穹之冠的银色地砖上，比任何时候都沉重——也比任何时候都安静。}


\textbf{"我跟你说过有三个名字是红冰刻的，"他对瓦尔德林说。"我一直在想——我为什么忘了他们。不是忘记——是被从历史中抹去了。如果他们被抹去了——今天所有这些站在这里的神——也可能会被抹去。不是被那个黑雾抹去——是被时间本身。被所有人忘记。而我——"}


\textbf{他没有说完。但他的父亲代他补完了："你不想让任何人再被抹去。"}


\textbf{然后是泽法利昂。他释放了他最后一发天穹之雷。不是射向敌人——是射向天空。那束雷光在裂霆高地上空炸开——不是爆炸，是绽放。像一个从未有人见过的、以雷为瓣、以电为蕊的巨花。它持续了整整一百息。一百息之后，雷光消退。天空变成了深沉的暗紫色——不是天色暗了，是泽法利昂将他剩余的所有天霆储存在这最后一百息的光芒中释放出来后，天空回到了它没有被雷光照亮之前的自然状态。泽法利昂站在消失的雷光下——他老了。不是皱纹老——是眼睛老。他的眼睛在那一百息间变成了淡灰色。}


\textbf{然后是达格德里安。他没有说任何话。他把右手按在大殿的地面上。从他的掌心——长出了一种没有土壤也能生长的根。那是幽翠圣林最后一棵圣橡的幼苗——它的母树在达格德里安离开圣林时就已经枯萎了。它把所有的生命力都灌入了这棵幼苗中。幼苗从达格德里安的掌心钻出来，发出极其微弱的、只有最敏锐的灵视者才能察觉的绿光。它还没有叶子。但达格德里安蹲下来——他高得多的身形弯下——他在这棵没有叶子的小苗边，闭上眼睛。在那一刻——如果有人同时看到他在幽翠圣林时那片他照看了无数个凌晨的圣橡林地——他会看到：达格德里安对这个大殿中所有人说的最后一句话——是用根语说的。不是用语言——是用所有根系能传递的、最深的那种振动。内容很短——只有四个音节。在翠庭的古老语言中，那四个音节不能被直接翻译为任何词汇——它更接近于一个系列：再见。谢谢。记得浇水。我不后悔。}


\textbf{莫丽甘娜没有自己走进帷幕。她做了一件更令人难以忘怀的事。她的所有渡鸦——翠庭领空内每一只曾经被她的灵焰触碰过的渡鸦——在这一天同时起飞。它们从翠庭的每一个角落飞向裂霆高地。它们在天穹之冠上空盘旋。数量之多，遮蔽了天空。在夕阳的最后一道红光即将触到地平线的那一瞬间——所有的渡鸦同时转向了极北之北。它们飞走了。不是消失在远方——是直接飞向那个正在逼近的黑色存在。莫丽甘娜在用她所有的渡鸦——和她自己——作为帷幕编织的"针"。当帷幕需要被从一端拉到另一端时，需要有什么作为针——什么快速、精准、不怕迷失、能够在两个世界之间穿越的东西。渡鸦是最合适的。莫丽甘娜走进了自己的渡鸦群中。在最后一刻——在渡鸦群即将消失在帷幕材料中的前一息——她转向了大殿中所有站在那里看着她的神。她笑了。}


\textbf{她说的最后一句话是："我在另一个通道那边等。你们慢慢来。"}


\textbf{然后是哈德昂——做了一个只有他自己能做、而其他任何人都无法做的决定。他不能在帷幕编织过程中停留在生者世界——超过一盏茶的时间。无光之渊的"无光"对他的灵焰是不可或缺的——离开太久他就会消散。但他还是来了。他用了他一整年积攒的、极其有限的能够在地面上停留的时间——站在天穹之冠的大殿中，献出了他身上所有可以献出的东西：他将无光之渊中储存的、他亲自收集了一万九千年的死者的灵焰印记——每一个死者在进入无光之渊时留下的那个微小的、只包含死者名字和最后遗言的灵痕——全部灌入了帷幕的基层。一共是六百亿零三千七百个印记。六百多亿个名字。六百多亿个遗言。它们被哈德昂一个一个亲手排列，一个都没漏。有的名字只有两个音节，有的遗言只有一个字——"暖"——那是一个在暴风雪中冻死的樵夫，他在意识消失前最后体感到的温度来自他妻子用来给他取暖的、在他死后三分钟内就也冻僵了的怀抱。}


\textbf{他的兄弟波里冬无法到达地面。他的灵焰被那枚远古符文"永系于水"囚禁在深海中。但他做了他唯一能做的事：在沉没之殿的最深处——在他永远不能离开的那个海底宫殿中——他将自己的灵焰中所含的、那些被深水冲刷了无数个世纪、变得光滑如镜的宁静感——灌入了包裹着帷幕之域的那片海洋。当诸神在裂霆高地上编织帷幕时，他们感到了脚下传来了一股微弱的、从海洋深处升起的低频震动。那是波里冬在用他的三叉戟击打沉没之殿的石砖——以他一直以来那种安静的、有韵律的、将孤独转化为力量的方式——为帷幕送上了最后一层由"水之记忆"塑造的涂层。}


\textbf{朱沃纳斯是最后一个进入帷幕的主要神祇。在他进入之前，他将《万域通则》交给了他能找到的最年轻的、还没决定是否要进入帷幕的小神——一位来自白石圣座最偏远边哨的守门少年，看起来不超过十四岁（按神的年龄换算——实际年龄大约是一千四百岁）。}


\textbf{"通则中有很多条款已经不再适用了，"朱沃纳斯说。"但我们不需要法律来继续存在。我们需要的是——"他指了指自己胸口——"这里的东西。把通则传给最需要它的人。不是最有权力的。是最需要的。你知道我在说什么。"}


\textbf{少年接过了厚厚的法典。他不知道他自己也不会活到能将法典传下去的那一天。但他在接过的瞬间——用他当时还没有意识到的方式——做了一个决定：他要将万域通则中的核心条款——那些最本质的、关于公平、关于宽容、关于不应该因仇恨而伤害无辜者的条款——记在他自己的灵焰中。这样即使羊皮纸在裂隙中朽烂，那些条款仍然有一个地方可找。}


\textbf{雅努维翁没有和其他诸神一起进入帷幕。他在最后一刻——在帷幕即将最终闭合的那一瞬间——做了只有他才会做的事。他走进了一道他自己建造的门。那是一扇没有标在万域通则任何条款中的、朱沃纳斯唯一没有对他提出过反对意见的门。}


\textbf{那扇门通向哪里？通向他亲眼所见的、在那一刻尚未被黑雾侵蚀的、帷幕之域的最后一小片净土。他把门开在了那片净土和帷幕的交界处。门的一面是帷幕之域。门的另一面是虚空。如果未来某一天有人需要逃出帷幕——或者有人需要走进帷幕去修补它——门在这里。它隐形的。它不需要锁。它只由雅努维翁自己的名字和意志守护。而他已经不再有意志了——他已经变成了门。你就是门。你在敲你自己。}


\textbf{所有六座圣座——霜脊、裂霆、翠庭、金沙、白石、雾霭——所有还在的、还愿意的、还舍不得这个世界的——都走了进去。}


\textbf{编织持续了九天九夜。在第十天的黄昏——当最后一丝神力的痕迹在大气中消散——帷幕闭合了。}


\textbf{帷幕之外的黑色存在——那些饥饿的、不可理喻的、尚未完全拥有的东西——被挡在了外面。}


\textbf{帷幕之内，是残存的土地。残存的凡人。残存的一切。}


\textbf{从此，诸神不再行走于大地之上。}


\textbf{他们变成了大地本身。}


\bigskip\hrule\bigskip


\chapter{卷二：帷幕之域}


\section*{第一章：新世界的黎明}


\textbf{帷幕闭合之后的第一千年，被称为"寂静纪元"。}


\textbf{没有神迹。祈祷不再被应答。霜脊山脉的玄冰巨城空了——那些曾经在塔尖上发光的冰晶已经失去了神力来源，变成了普通的冰。裂霆高地的银色圆顶不再被闪电环绕——泽法利昂留下的最后一道雷在两百年后也消散了。幽翠圣林的圣橡已经全部枯死——达格德里安将它们的根系也编织进了帷幕，所以它们的树干仍然站立在那里，但它们是空的——在灵质层面中，它们只是一些被形式困住的记忆。}


\textbf{金沙圣座的金色沙粒——那由亿万个被遗忘的名字构成的沙地——风化了。名字一个接一个地从沙粒表面脱落，飘入帷幕中。今天你如果走过骨灰平原——金沙圣座的旧地——你偶尔可以看到一撮细沙在你脚边升起一小股风。那是一个名字在离开。它是一个曾经被某人记住、然后被遗忘了几千年、终于决定要放弃成为名字的——名字。安息吧。对不起。没有人记得你了。}


\textbf{白石圣座的白色大理石宫殿——朱沃纳斯的律法之城——在寂静纪元的前两百年中还维持着原样。但大理石的灵质记忆是有限的——当刻石的神力消散，石头恢复成了普通的石头。风化开始。最初的五百年中，宫殿的正门塌了。之后的五百年——穹顶裂了。一千年后，整个白石圣座变成了废墟。但有一个地方保留了下来——在废墟深处，一扇被厚厚灰尘遮蔽的木门。门的材料不是大理石的——是一种没有人能辨识的木材。它上面用只能在水下看到的墨水写着一个字。不是字。是字的痕迹。是那个被朱沃纳斯本人刻在万域通则最后一页的、被故意留白的字。这扇门至今仍然可以打开。但它通向哪里——没有人知道。所有打开过的人，都在门的另一边找到了一个他们以为此生不可能再见到的人。}


\textbf{雾霭列岛的仙岛在帷幕闭合后沉没了。不是被淹没——是主动下沉。玉穹尊者在献祭前给雾霭列岛下了一个"沉潜咒"——所有的仙岛会在大约三百年后沉入海面之下。不是消失——是等待。当某一天——玉穹尊者相信会有那一天——有人找到了唤醒雾霭诸神的方法，仙岛会重新浮出水面。届时，那些在海底沉睡了不知多少年的石径会重新连成完整的步道。届时，雾中会有一个人——站在第一个台阶上——对你说："你来得正好。茶刚泡上。"}


\textbf{那一天至今还没有到。但玉穹尊者的茶——在最深的海底的一个密封的玉瓶中——还在冒着热气。}


\textbf{这就是帷幕闭合后的世界。诸神已去。他们曾经守护的土地现在由凡人自己来守护。}


\textbf{没有人问过凡人是否愿意。没有人给过他们选择。}


\section*{第二章：破碎之地的地理}


\textbf{帷幕之域是六座圣座旧地碎片的拼接体。就像巨人将不同颜色的黏土揉在一起、然后随意摊开——每一个区域都来自不同的圣座、有不同的地貌、气候和灵质残留密度。}


**\textbf{永冬山脉}**\textbf{在最北端——这是霜脊圣座的核心遗址。山脉高耸入云，峰顶终年积雪。除了风声和偶尔出现的、被冰冻在冰层中的远古巨兽骨骸之外，这里几乎没有任何生命的迹象。山脉深处的某个冰洞中——据说到今天还保存着瓦尔德林在献祭之前留下的最后一块霜纹石板。找到它的人，可以得到瓦尔德林最后一条未说出的智慧。但那个冰洞的位置时刻在变动——不是物理上的变动，而是一种只有瓦尔德林的灵焰残响才能解开的、基于月相和极光角度的位置加密。这块石板至今没有被找到。}


**\textbf{风嚎荒原}**\textbf{是霜脊山脉和裂霆高地之间的过渡地带。这里的气候极其恶劣——狂风终年不断，风速大到能将岩石表面切割成光滑的镜面。荒原上偶尔能看到被风蚀过的、属于某个无人记得的文明的石柱。石柱上的文字早已被风磨平。但如果你在特定的日出时分——当太阳的第一道光以特定的角度——不，不是角度——是温度——照射石柱时——石柱的阴影会在地面上短暂显现出几个字。那些字不在石头上。它们在影子里。用影子写字——这是风嚎荒原上那支早已灭绝的文明的语言。没有人能解读。但每当影子显现时——那一瞬间——你会觉得有人在你身后。不是威胁你的人。是一个你很久以前认识的、以为这辈子不会再见到的人。}


**\textbf{裂霆高地}**\textbf{至今仍然是最壮观的遗迹之一。泽法利昂的天穹之冠的银色圆顶虽然不再发光，但它的金属外壳仍然完好无损——因为那种银色是一种带有灵质记忆的金属，它"记得"泽法利昂的灵焰频率，所以一万多年来一直没有氧化。高地的各个裂谷中，据说到今天还能找到泽法利昂投出雷矛后在岩石上留下的焦痕。焦痕的形状千奇百怪——有一道焦痕恰好形成了一个完美的人形轮廓（泽法利昂的追随者说那是他曾经救过的一个凡人的身影），另一道焦痕的形状和霜脊山脉南麓的地图一模一样（他在某夜的雷暴中想家了）。}


**\textbf{无光之渊}**\textbf{——哈德昂的旧领地——的入口已经在大地震中崩塌了。但有几个被称作"渊息孔"的微型裂隙留在地面上——每隔几十里就有一个——从里面断断续续地冒出微弱的、携带着已死之人的低语的空气。站在渊息孔旁边，你可以闭上眼睛，静听。如果幸运——或者说，如果不幸——你可能会听到一个你认识的人的声音。那不是鬼魂——哈德昂说过，无光之渊中没有鬼魂，只有灵焰残响。那些残响中的信息每次被"渊息孔"释放出来时，只包含一段信息：死者在死之前最后一个完整的念头。有的人在死之前想的是"忘了锁门"。有的人是"我爱"。还有的人——只有一声叹息。不是悲伤的叹息。是那种——"原来是这样"——的叹息。}


**\textbf{沉没之殿}**\textbf{——波里冬的老家——现在位于一片叫"无尽海"的巨大水域的深处。海面上的水是普通的海水——但如果你潜入水下大约三里，你会发现海水开始变了味道。不是咸的——是什么味道都没有的。那是最纯净的水——波里冬死前从沉没之殿中排出的、被他用最后的神力净化的"初水"。据说喝了初水的人能恢复灵焰中最核心的那一点纯净——不管他已经被蚀痕侵蚀了多大面积。但沉没之殿在三里之下——没有人能潜到那里还活着回来。波里冬的孤独——在他生前无法被分担——在他死后也无法被触及。}


**\textbf{幽翠圣林}**\textbf{的残骸现在被称为黑森林。名字不准确——树不是黑的。它们是深灰的、空的、但依然站着的。圣橡虽然失去了灵质层面的枝叶，但它们的物理树干没有倒。一万多年来，它们的树皮从棕色逐渐变成了银灰色。在正午时分——当阳光恰好穿透林冠——银灰色的圣橡树干上会短暂地浮现出根语留下的最后一道信息。那是一幅图——不能用语言描述——是所有圣橡在被编织进帷幕前的最后一瞬间，对彼此同时"说"的一句话。那幅图的内容——据所有见过的人说——是一个圆圈。不是几何意义上的圆。是"圆满"的圆。是"我们曾经一起活过、一起死过、不后悔"的圆。}


\textbf{黑森林深处有一个地方被称为"莫丽甘娜的渡鸦冢"。那是翠庭所有渡鸦在献祭之后——它们大部分都死去了——的尸体被风干后堆积而成的一座小山。山的形状——精确地形成了一只俯冲的渡鸦的轮廓。据说在血月之夜，冢上的渡鸦羽毛会发出一种暗紫色的微光。那不是任何灵质的残留——那只是莫丽甘娜在献祭前给她的渡鸦群的最后一道命令，一道不需要神力就能执行的、极其简单的命令："记住我。"}


**\textbf{骨灰平原}**\textbf{是帷幕之域最令人不安的地方。它曾经是金沙圣座的旧地，但在诸神黄昏中，金沙圣座的核心——那口存放着所有记录的石室——从地下被翻到了地表。石室破碎了。所有的白石板上刻着的灵符——那些记录了金沙两万年历史的符号——混入了沙土中。如今骨灰平原上铺着细密的灰白色沙土——那不是沙，是骨灰。是亿万个名字风化后剩下的最后一点粉末。在骨灰平原上行走时，脚下偶尔会踩到硬物——那是残存的灵符碎片。如果你捡起碎片，对着月光看——你可能会在碎片表面看到几个微微发光的符号。伊希丝的灵符是设计用来在黑暗中自然发光的——所以一万八千年后，它们还在发着极其微弱的光。你能读到的内容取决于你捡到的是哪一块碎片——可能是某一年谷物的收成，可能是某一位早已被遗忘的法老的婚礼日期，也可能是一个母亲的笔迹，记录她两岁孩子的第一个完整句子。}


\textbf{骨灰平原上最重要的禁忌是：不要在平原上久留。不是迷信——是物理（灵质物理）上的。骨灰——那由无数被遗忘的名字构成的粉末——会缓慢地吸附活人的灵焰。不是噬灵式的强行剥离——是大自然的缓慢温柔的不动声色的吸取。你待得越久，灵焰流失得越多。第一天你不会察觉——你的灵焰减少了不到万分之一。第五天——你开始觉得精力不济。第十天——你的眼睛会看到颜色比之前淡了一层。第三十天——你会变成一具还在呼吸的、但灵焰已经薄到只剩一层蝉翼的空壳。}


**\textbf{白石帝国遗址}**\textbf{——朱沃纳斯的旧城——是帷幕之域中保存得最好的废墟。不是因为有什么神力保护——是因为朱沃纳斯生前用的建筑材料太结实了。那种白色大理石——"万律石"——是他亲手从星空中坠落的一颗小行星上采来的。白石帝国遗址最著名的地标是"律法广场"——一片曾经容纳过上万人同时辩论的巨型石面广场。广场上的讲坛——朱沃纳斯曾经站过的地方——至今还留着他最后一次演讲时的脚位印记。两个浅坑——左脚的浅，右脚的更浅（因为他演讲时重心微微向前倾，像一个即将起飞但还在地面上的人）。浅坑中的石面比其他地方光滑——不是风化的——是无数次有人跪在那里、双手按住他的脚印、问他"我们现在该怎么办"而磨光的。很多人来过。没有人听到过回答。他们都说从来没指望会得到回答——去那里跪不是因为期待答复，而是因为这是一个曾经有神存在过而那个神愿意听凡人说话的最后的坐标。}


**\textbf{雾隐禅院}**\textbf{是雾霭圣座下沉后唯一露出海面的建筑。它建在一座微型礁岛上——岛上唯一的建筑是一间木造禅房和一个已经长满苔藓的石灯笼。禅院中曾经住着玉穹尊者最后一任弟子——一位在玉穹尊者献祭前三天才正式入门、还没有来得及学到任何法术、只被传授了一种呼吸法的年轻僧人。没有人知道他叫什么。他自称"虚舟的师父的师父"——但他自己的名字，他说"不重要，因为我已经不能教任何人了"。他每天做的事是清扫禅院前的石径。那条石径——当雾霭列岛还在海面上时——连接着所有的仙岛。如今石径的台阶大多在水面下，只有禅院所在的这节还露出水面。虚舟的师祖每天早上把石径扫干净——不是为了有客人来（一千年没有客人了），而是因为扫石径这个动作——那种规律的、不需要思考的、将落叶从石面上扫去的机械动作——是唯一能让他平静下来不去想玉穹尊者最后说的那句话的事情。}


\textbf{那句话是什么？年轻的僧人从来没有告诉过任何人。他只是在每次扫完石径后，坐在石灯笼边，给自己泡一壶最淡的茶。茶淡到几乎就是白水。但他说玉穹尊者最后告诉他的那句呼吸法——"呼入时想象你正在吸入世界上所有的安静；呼出时想象你在呼出你所有的不安"——只有配这壶最淡的茶才能发挥最大效果。他活到了一百三十岁。临死前他写了一个字条——压在石灯笼下。字条上写着："茶还热着。进来吧。"}


\textbf{没有人知道他是写给谁的。但石灯笼里的蜡烛——直到今天——仍在燃烧。不是神力。是那些来到禅院的旅人——在一千年里，大约有几百个——每个人都给蜡烛加了一点自己的融蜡。一千年。一个人不会烧完一根蜡烛。一千个人每个人贡献一点——蜡就永远够用。这就叫"永恒"——不是不会被用完，是每一个路过的人都知道要给下一个人留一点光。}


\textbf{所有这些区域——北从永冬山脉到南至骨灰平原，东从雾隐禅院到西至黑森林——都围绕着一个共同的中心。}


\textbf{在那片盆地中的聚落——叫做暮色聚落。}


\section*{第三章：暮色聚落}


\textbf{暮色聚落不大。快走的话，你可以在不足一刻钟内从它的一头走到另一头。}


\textbf{十二栋屋子围成不规则的一圈。聚落中心是一棵早已枯死的巨大黑色树干。那不是普通的树——那是世界树残骸上掉落的一根枝桠，被第一代聚落居民插在土里。所有人当时都认为它不可能活。它也确实没有活——没有根，没有叶，没有任何属于"活着的树"的特征。但它在每年春天都会在枝顶冒出几朵纯白色的花——不多，三四朵。花期极短，三天就落。落地时花瓣铺在黑色树干周围的泥土上，薄薄一层白。那不是雪，但那看起来像雪。而在这个一年之中十一个月都是暗红色月光的聚落中——一抹白色是一次比任何愈灵汤剂都更有效的心灵疗愈。}


\textbf{聚落外面是荒原。荒原上有几条被无数夜行者踩出的、浅浅的路径。这些路径通向各个灵质热点：水井、铁匠铺、圣橡遗址、观测塔、墓地。其中两条——一条向北、一条向南——分别通向两个最重要的目的地。向北通向聚落边缘的"观测点"——奈莉娅的观测塔。向南通向"放逐台"——一个被蚀刻了古老符文的石台。当聚落公投决定将某人放逐时，被放逐者被押送到这里。不是被杀——是被"放逐"——即被从聚落的集体灵焰感知网中永久断开。帷幕会感应到这个被断开的人——它会认为这个人的灵焰不属于聚落——于是不再保护它。夜幕降临时，裂隙中的东西可以自由地捞走这具无人认领的灵焰。}


\textbf{聚落外有一座孤独的哨塔。那是帷幕守卫们值夜的地方。哨塔的石砖上刻满了之前历代守卫的名字和遗言——有的是用刀刻的，有的是用自己灵焰烧的，有的只是用炭笔刷刷写的——炭迹已经被风吹淡了，但名字还在。你如果用手摸那个名字，可能会在指尖感受到一阵微弱的振动——那是名字的主人在刻字时留下的灵焰印记。最快的解读速度——帷幕学者大约需要小半个时辰——可以把印记中包含的最长的一段信息译出来。一般那都是一些简单的、短的、被刻在石头上因为刻字者觉得"如果我明天死了至少有人知道我想过什么"的话。}


\textbf{聚落的居民们世世代代就这样生活着。种地，打铁，织布，烤面包，讲故事，守夜，等天亮。他们都知道世界已经碎了。他们都知道诸神已经死了。他们都知道——当夜幕降临——有些可怕的事情会发生在一些人身上。但他们仍然起来了——每一个清晨在血月的光变得能被直视之后——继续生活。不是因为勇敢。是因为没有别的选择。是因为"还有人在等你回家"是一句在这个聚落中比其他任何地方都更真实的陈述。}


\bigskip\hrule\bigskip


\chapter{卷三：灵焰的发现}


\section*{第一章：目盲者所见的火}


\textbf{在帷幕闭合后的约第四个世纪，一个改变了人类对自身理解的事件发生了。}


\textbf{一个目盲的女人——她走路时需要用手杖叩击地面、需要一只老狗牵着绳为她引路——在某个血月之夜后在聚落中说了一句让所有人都愣住的描述。}


\textbf{她说她能"看见"一种光。}


\textbf{不是用眼睛——她没有眼睛。她出生的那年，一种叫"幕咳"的疾病（一种由帷幕波动引起的、在今天早已被草药学者们消灭的疫病）导致她尚未发育完全的眼睛在子宫中停止了生长。她的眼眶是空的——眼睑覆盖着平滑的皮肤，下面什么都没有。}


\textbf{但她坚持说她看得见。不在白天——白天她什么都看不见。是在夜晚——在月亮开始泛起血红色的光、帷幕最亮的那个时段——她能"看见"。不需要眼睛。不需要光。不需要任何人告诉她在看什么方向，她就能准确地指向聚落中的每一个人——"你的光——是金子色的。你的——是淡蓝的。那个人的——"她停住了——"是黑的。"}


\textbf{她说她能在每一个人体内看到一团"火"。不在胸中——在更深的地方。在某种比骨骼更内层、比血液更隐秘、比呼吸更原始的层次。她说那团火是一切的中心——它比心脏更根本，比大脑更古老。它是生命的源头——"它先于你的第一次呼吸就存在。它在你死后还会延续一小段时间。它不是你的——你是它的。你只是它借以存在于这个世界的一件外衣。"}


\textbf{她最后下的结论——十多年后得到了帷幕学者奈莉娅的实证——是：一个人的灵焰（她甚至自己也起了这个名字——因为她说它看起来"像火焰，但不是热的，是在燃烧一个东西但从不烧完的那一小片火焰"）的颜色、亮度和形态反映着这个人的本质。纯净者——灵焰清澈的、未被任何外来伤害者——灵焰是淡金色或亮白色的。被某种负面经历伤害过的人——灵焰偏灰。濒死之人——灵焰极弱，已经看不清楚边界在哪。而被一个她从未见过但能清晰感知到的东西感染的人——灵焰边缘被一圈黑色包围。}


\textbf{那个黑色——她跪在那些灵焰变黑的人面前，把没有眼睛的、空洞的眼窝朝向他们——"不是你们的。是外面的。是从外面爬进来的。我能闻到——它闻起来像烧焦的蜂蜜。"}


\textbf{目盲的艾瑟（她的名字被后人神圣化——"艾瑟"在当地古语中意为"先见者"）在她二十四岁时死去了。但她的发现被保留了下来。此后一代又一代的灵视者们继承了她的能力——每个人看到灵焰的方式都不同，有的看到颜色，有的听到声音，有的通过触碰感到温度，有的在睡梦中梦到别人的灵焰形状——但所有人都确认了一件事：那团火是真实存在的。它叫做}**\textbf{灵焰}**\textbf{。}


\section*{第二章：灵焰不倒}


\textbf{奈莉娅——第一位系统化灵焰理论的帷幕学者——在艾瑟死后约四百年出生。她花了她人生的头三十年观测和记录灵焰的各种性质。她的观测塔——旧霜脊圣座的一个瞭望塔废墟——现在成为了帷幕之域灵焰研究的总中心。}


\textbf{奈莉娅最重要的发现——后来被称为}**\textbf{奈莉娅第一定律}**\textbf{——是：灵焰不可被创造，不可被毁灭。当一个人死亡时，他或她的灵焰不会真的"消失"。它会离开躯体——一段短暂的悬浮后——要么被帷幕接收（升入帷幕中，与诸神的残留意念融合），要么被裂隙中的东西趁机捞走（坠入混沌中，被非存在的力量当作"存在样本"来吞噬）。}


\textbf{"没有人真的'死去'——如果你把死亡定义为存在从世界中被彻底删除，"奈莉娅在《观测录》第一卷中写道。"死亡只是灵焰更换了住所。正如一个人从一间漏雨的旧屋搬到另一间可能更好、也可能更差的新居所。而我们——这些还在旧居中的人——应该做的事不是哀悼搬家的人，而是为他新居的屋檐祈祷晴好。"}


\textbf{这个理论——灵焰不灭——给聚落带来了巨大的安慰。但它也带来了同样巨大的恐惧：如果灵焰不灭，那么被裂隙中的东西捞走的灵焰——那些被黑雾吞噬的灵焰——并没有消失。它们去了那里——在某种我们无法感知的状态下——继续存续。是在怎样的条件下存续？它们能感知到自己的存在吗？它们在痛苦吗？没有人知道答案。但这个问题——"那些被捞走的灵焰在那边是什么感觉"——让很多守幕者宁可献祭自己于帷幕，也不愿被噬灵。}


\section*{第三章：蚀痕的出现}


\textbf{蚀痕不是外来的物理侵入。它不来自任何已知的病原、毒素或创伤。它是一种灵质层面的感染——由裂隙力量引起、通过灵焰接触传播、最终导致整个灵焰被一层不断扩散的黑色物质覆盖。这层黑色物质就是"蚀痕"。当蚀痕的扩散到达灵焰的中心——那个叫做"灵核"的微小的、极热的、淡金色的点——这个人的灵焰就发生了不可逆转的质变。从内而外——灵焰从金色变成全黑。这个人不再有自己的灵焰。取代灵焰的——是一个由蚀痕构成、由一种寄生在蚀痕中的微小碎片驱动的"反灵焰"。}


\textbf{这个反灵焰不再向外辐射热量和光明——而是向内吸收。吸收什么？吸收周围所有活物的灵焰。它像一口倒反过来的深井——不是从底向上涌水，而是从顶向下吸入。吸入的速度取决于这个蚀变者与目标之间的灵焰频率的相似度——越相似，吸力越强。所以蚀者更倾向于攻击那些与他生前有感情纽带的人——家人、密友、恋人——因为他们的灵焰频率与他本身的（在他被蚀变之前的那个频率）是极其相似的。}


\textbf{蚀痕的传染机制至今是帷幕之地最重要的未解之迷。但奈莉娅的观测提供了关键的线索：蚀痕不是主动入侵的——它是被许可的。没有一个人的灵焰是在未经历某种形态的"灵质破裂"的情况下就被蚀痕附着的。灵质破裂——奈莉娅定义为"灵焰在遭受极端情感冲击后的暂时性裂纹"——为蚀痕提供了附着的缺口。最深的灵质破裂发生在那些经历了深度丧失或深度背叛的人身上：失去至亲、被最爱的人欺骗、在最绝望的时候被抛弃。这些经历在灵焰表面撕出了微小的、肉眼不可见的裂隙。而那些在帷幕之外等待的混沌力量——它们能在极远处感知到一个灵焰上的每一道这样的微小裂隙——它们会派出蚀痕碎片进入裂隙中，像孢子落入沃土。}


\textbf{这就是为什么有些人被蚀者反复攻击却从未被感染——他们的灵焰太完整了，没有裂隙可以供蚀痕攀附。这也是为什么最严重的蚀痕感染往往不是发生在那些与蚀者正面冲突的人身上——而是发生在那些在情感或道德层面遭受了重大打击的人身上。背叛——比任何噬灵攻击都更有效地撕开灵焰。被朋友放逐、被恋人的谎言击中、在绝望中选择出卖别人——这些灵质层面的自毁行为——每一次都在喂给帷幕之外的存在一个全新的、温暖的缺口。}


**\textbf{第一个被记录的蚀变者}**\textbf{是一个铁匠。他的名字——按照帷幕之地的传统，蚀变者的名字会被从官方记录中永久移除，因为"名字本身是灵焰的一个链路——如果名字还在流传，那个蚀变者就能通过名字来感知活在记忆中的人"——已经被烧毁了。仅存的代号是"第一人"。他经历了什么我们已不可知——他的名字、他的故事、他蚀变的原因都随他的名字被烧毁而永远消失了。但他在蚀变后写下的那两个字——"谢谢"——至今仍在奈莉娅观测塔最底层的一间封存室中被单独保存在一个密封的水晶盒里。打开盒子的人——一个都没有。不是不敢。是不知道自己如果打开——看到了写在石头上的两个字——会做出什么反应。奈莉娅说她也不知道。}


\bigskip\hrule\bigskip


\chapter{卷四：蚀月纪元}


\section*{第一章：噬灵的第一次记录}


\textbf{聚落历史第五百五十七年——这个数字如此精确是因为那之后每一年都被刻在了观测塔的记录石板上——发生了一件至今被老人们在炉火边以最低沉的嗓音重述的事件。}


\textbf{七个蚀变者在同一个血月之夜走出了他们的屋子。不同的屋子，不同的时刻，互不通信——但他们都在同一个夜晚做出了同样的事。他们以某种不是行走的方式穿越了庇护所的墙壁——不是因为墙壁不坚固，而是因为他们的灵焰被蚀痕完全改造后不再与物理实体碰撞。对于蚀者来说，一堵石墙和一蓬雾是一样的——都是"不是灵焰"的东西，都可以被穿过。}


\textbf{七个目标。七个不同的清晨发现的尸体。七个人——毫无外伤、毫无挣扎痕迹——在各自的庇护所内失踪了。只留下他们的身体躺在原处。身体是完整的，心脏仍在（不是物理意义上被取走），但灵焰已不在。他们的灵焰——被强行从身体中剥离——成为了蚀者灵焰中的新碎片。碎片悬浮在蚀者的灵焰中——不融合，不消化——像被投入深水中的一滴油。它持续地、微弱地发出原主人最后那个瞬间的记忆回响。蚀者不会解读这些回响——但帷幕之外的存在会。}


\textbf{这就是}**\textbf{噬灵}**\textbf{——后来被奈莉娅定义为"灵焰的强制剥离与吞噬"——首次被记录的事件。}


\section*{第二章：守幕者的誓言}


\textbf{七个死者的家人和朋友在之后的几个月中陷入了比哀伤更黑暗的状态——不是在哀悼死者的离去，而是在恐惧自己是否是下一个。}


\textbf{聚落中的灵视者们——到此时已经约有十几人——提出了一个集体的想法：将灵视者组织起来，建立一个专门对抗蚀痕的团体。不是军队——帷幕之地没有军队，也不需要军队。蚀者不能用军队来打——因为蚀者不占据领土、不争夺资源、不战斗。他们只是在夜深人静的时候——穿过你的墙壁——取走你的灵焰。你一个人怎么用军队？你不需要军队。你需要的是能看见他们、能在灵质层面找到他们、能将他们挡在外面——在不可能的情况下将你从他们嘴里抢出来的——专门的人。}


\textbf{八个人成为了守幕者的最初成员。他们彼此之间的信任——在最开始——仅限于能一起待一个晚上而不互相检举对方是蚀者。}


\textbf{第一位成员是那个叫卡珊德拉的女孩。她在十二岁时继承了目盲艾瑟的灵视天赋——能在没有光的环境下看见灵焰的颜色、形状、流动。到了二十六岁——在观测了上万个灵焰样本之后——她能将一个蚀痕的"边界"精确到一根发丝的宽度。她能在灵焰的一个边缘上分辨出哪些黑色是蚀痕（需要警惕）哪些深色只是那个人的灵焰天然的颜色（雪松灰、檀木褐——每个人的灵焰颜色不同是由他们的灵质祖先决定的）。她的老师奈莉娅写道：卡珊德拉的左眼在一次观察中因为过于靠近一个活跃裂隙而导致瞳孔永久僵化——此后她的左眼瞳孔不能缩小或放大。但它获得了另一种功能：它在特定角度下能看到帷幕表面上正在被动阅读的蚀痕碎片——像一个人在书的背面看到透明的字迹。}


\textbf{第二位是}**\textbf{草药学者埃洛伊丝}**\textbf{。她的师祖曾在霜脊圣座旧地的冰裂缝中发现了一种只生长在零下八百度的、靠微弱神血残留维生的灵植——"霜星苔"。霜星苔被移植到暮色聚落后，埃洛伊丝的师祖花了五代人的时间繁育它。埃洛伊丝本人发现了霜星苔的花粉当施撒在一个蚀者的灵焰上时——花粉中残存的神力残留（来自瓦尔德林最后的霜纹碎片）会与蚀者灵焰中的蚀痕发生一种猛烈到瞬间即毕的反应。蚀者灵焰中所有的蚀痕在几息之内全部烧毁。但这个反应有一个悲剧性的副产品：它同时也烧毁了蚀者灵焰中原主人的残存意识——那些还在蚀痕深处被困着的、最后的、再也说不出任何一个字的灵焰原主人——被永远释放了——但不是回到帷幕，而是化为灰烬。埃洛伊丝在笔记中写道：她每次使用蚀灭药剂时，她都要对被毁灭的蚀者说一声"恕我"。即使那个人已经无法听到。即使她自己也知道——她必须做这件事。有些哀悼是留给敌人身上那最后一点未曾堕落的残余。}


\textbf{第三位是}**\textbf{愈灵师奥尔德温}**\textbf{。他能在灵质层面触碰别人的灵焰——这是一种与生俱来的天赋，不是后天训练出来的。他的触碰不只是诊断性的——他可以通过将自己的灵焰与对方短暂联结来修补对方灵焰上的裂缝。原理类似于一个温手暖一只冻手——奥尔德温的灵焰本身纯净，它的纯净度会对受伤的灵焰形成一种"参照压力"：受伤的灵焰感知到另一团更纯净的灵焰在它旁边，就顺着那个参照来修复自己。但每次他做这件事，他的灵焰必须在对方灵焰中待至少三息——在这三息内，他自己的灵焰暴露在对方的灵焰状态中。如果对方是蚀者——即使蚀痕只有浅浅一层——奥尔德温在三息内会被感染。他一生大约治愈了上百人。他死于七十岁——不是死于蚀变，而是老年的灵焰自然熄灭。这是他应得的。他治愈的人中有三分之一成为了新的愈灵师——把他传下来的手艺接了下去。}


\textbf{第四位是}**\textbf{帷幕守卫马尔库斯}**\textbf{。他能将自身灵焰延伸出体外形成一个薄到透明的保护罩——最多罩住一个人。保护罩能挡住一次噬灵攻击——不是反弹，是吸收：吸收的那片灵焰被撕裂的能力转移到马尔库斯自己的灵焰上。每一次成功的守护，都在他的灵焰上留下了一道新的裂缝。这不是蚀痕造成的裂缝——裂缝本身不带有蚀痕，它只是一种结构疲劳。当疲劳积累到临界点——他的老师奥尔德温说——"你的灵焰会碎。不是被蚀——是碎。碎在你自己身上。"马尔库斯活到了四十八岁。他死于灵焰碎裂。死之前他在哨塔上守了一整夜——没有人知道他守护的是哪一个屋子。他不肯说。他只重复了一句话——在嘴唇从内部裂开、呼吸变得不稳定之前的最后一句："往北走的人需要南墙。"}


\textbf{第五位是}**\textbf{灵痕追猎者塞拉斯}**\textbf{。他从小就生活在黑森林——没有任何人类社会。他被圣橡遗留的灵质残响养大。他能看到灵痕——灵焰在空间中留下的残迹——一条发光的足迹，颜色从淡金（纯净者）到灰（刚刚感染的）到深黑（蚀者）。他将灵焰弩造了出来——材料是黑森林中一棵枯死圣橡的心材，加上他自己的灵焰残响。发射的不是箭——是将他自己灵焰的一个极小碎片在弩弦上凝成烙印、加速、投射。灵魂的弩。不是打身体——是打灵痕。一个蚀者在噬灵后、移动时留下的灵痕——如果被灵焰弩的弩箭射中——会从弩箭的击中点燃起连锁反应。不是烧他——是烧他的灵痕。蚀者丢失灵痕后——他就像人在黑暗中丢了拐杖一样——无法再追踪下一步。他会在原地兜圈——困惑地——直到早晨的日舟之光将他烧尽（蚀者不能暴露在日光下——拉赫卡里昂的日舟虽然因帷幕阻隔已经远不如旧日明亮，但其中残存的、来自两万年前金沙圣座的神力余热依然能将蚀者的反灵焰烧成灰）。}


\textbf{第六位是}**\textbf{灵织者}**\textbf{——那些拥有灵焰感知、帷幕低语、灵织术的人。他们构成了守幕者的大多数。每一个灵织者都有自己的感知特长：老兵灵织者能用"嗅觉"探出蚀痕；行路灵织者能在一夜之间两处来回收集信息；学徒灵织者能短暂稳住一个濒死者的灵焰；记述灵织者在聚落会议上能说双倍的时间；守夜灵织者能察觉到门外经过的人有多少；炊火灵织者通过灵焰烹饪的痕迹感知人们的去向；锻甲灵织者能打造暂时挡住蚀者的灵质门锁；织幕灵织者能将碎片信息编织成预测。灵织者们联合在一起——通过灵焰共鸣——形成了一个覆盖整个聚落的灵焰感知网络。没有这个网络，蚀者可以毫无障碍地在聚落中自由穿行。}


\textbf{最后两位——是蚀者和冥僧人。不是守幕者——是在另一边的。}


**\textbf{蚀者}**\textbf{不是组织。蚀者没有首领。他们是灵焰被完全蚀变后的产物——每一个蚀者都是独自蚀变的。他们之间唯一的共同感知是"裂隙共鸣"——在夜空中，当一个蚀者与另一个蚀者的灵焰在裂隙的同一频率范围内共振时，彼此能感觉到对方的存在。感觉不是精确的——更像是一种"有人在那边"的直觉。有时候两个蚀者在同一夜走进同一间屋子——他们会在门口对望。没有语言。没有竞争。一种纯本能层面的理解：这个目标够大，我们两个都可以咬一口。}


**\textbf{冥僧人}**\textbf{则完全不同。冥僧人在他们的残余人性碎片彻底耗尽之前，是蚀者中最清醒、最理性、最善于策略的一类。他们是主动做了交易的人——不是被蚀变的，是自己走上那条路的。当一个冥僧人的人性碎片耗尽、灵焰中的蚀痕到达了核心——他就不再是冥僧人了。他变成了一种比超级蚀者更难以对付的存在——一个"虚空容器"——他的身体不再是蚀者的身体，而是帷幕之外某个存在的临时躯体。那个存在——不管它是什么——通过这具躯体暂时地看见了帷幕之内的世界。}


\bigskip\hrule\bigskip


\chapter{卷五：八道源流}


\section*{第一章：察灵之术——帷幕学者的传承}


\textbf{帷幕学者的能力并不属于那些被"选中"的人——不是神灵的赏赐，不是命运的安排。它是可以被发现的。任何人在最深的安静中——在那种你几乎能听到自己的灵焰在体内燃烧的声音的安静中——尝试用灵视去感知另一个生命的灵焰，都可能在某一天突然"看见"。}


\textbf{但这不意味着每个人都能成为帷幕学者。就像每个人都能学会唱歌，但不是每个人都能唱到能让人落泪。帷幕学者的察灵之术——那个能"看到灵焰的蚀痕边界"的能力——需要几十年的训练。训练者首先要学会熄灭自己内心的噪音。一个焦躁的人看不准——因为自己的灵焰波动会在被观测对象的灵焰上投下"观测者自身频率"的干扰波纹。一旦干扰存在，观测结果就失真了。看岔了几根头发丝的蚀痕边界——可能把纯净者误判为蚀者——或者在更糟糕的情况下，把一个已经蚀变了七成的人当成"边缘灰色"来判定而无视了——而那个人的第一次噬灵就在当晚。这样的误判不会只发生一次。}


\textbf{所以帷幕学者必须先从"清心"开始——在观测塔中连续静坐十二个时辰（一个完整的日夜周期），直到所有杂念从灵焰表面沉到底部。然后开始"观灵"——从最容易的目标开始（一棵圣橡残枝——灵焰虽然微弱但稳定可预测）到最难的（一个即将蚀变、灵焰正在翻涌的临界者——探测他的灵焰就像在激流中测一瓢水的温度）。这个过程需要五到十五年，取决于个人天赋。卡珊德拉花了七年。奈莉娅花了十一年。有些人在第十七年的修炼中仍不能稳定地定位蚀痕边界——他们在第二十年准时离开了塔，加入了灵织者的行列。不是失败——是找到了一条更适合他们的路。守幕者中没有人会因为你不适合这个角色而被视为失败。你只是去了另一个岗位。}


\textbf{最顶尖的帷幕学者能做到的不只是察灵。他们还能"溯灵"——在灵焰中反向追溯一段被深埋的记忆。不是读取记忆——那已经超出了任何人（包括死者）的能力。而是"定位"——定位到灵焰中有一段被刻意封存的东西。它被封存的原因可能是太痛苦、太羞耻、太可怕——或者太重要。被封存的灵焰片段不会像正常灵焰一样"流动"——它凝固在灵焰中被封存的结节点，像一个被琥珀包裹着的微小的、不能动弹的飞虫。溯灵者能从上方——从一个安全的角度——看到那个结节点中朦胧的轮廓。好的溯灵者能从轮廓中分辨出那是什么类型的记忆：是丧失（淡蓝色的轮廓，边缘锯齿形），是被背叛（暗灰色的轮廓，带有尖锐的刺芒），是在最后一刻选择了善良（柔和的金色轮廓，周边像被一阵暖风拂过）。溯灵者从不告诉被溯灵的人他们看到了什么——不是秘密，而是因为那些被封存的东西被封存是有原因的。没有人有权利在没有被允许的情况下揭开别人的旧伤疤。溯灵者只被允许做一件事：在观测日志中注明"此人有封存物"或"此人灵焰纯净——没有明显的封存状"。如果当事人自己愿意谈，溯灵者可以陪着他，但不能替他开口。这是一个沉默的执业宣誓，所有帷幕学者在第一次溯灵前都会被要求念一遍："我看到的——是你还没有准备好说的。我不会替你说。"}


\section*{第二章：灵植之道——草药学者的传承}


\textbf{埃洛伊丝的师祖不是帷幕之地的本地人。她来自早已被裂隙吞噬的霜脊冰原——一个在霜脊山脉最深处还有人类居住的微小聚居区。她的家族在她之前四代就开始研究冰原上的特殊灵植。这些植物之所以能在没有阳光、没有液态水、没有普通土壤的极寒环境中存活——是因为它们不是靠光合作用活的。它们靠的是极其微弱的、已死诸神的残留意念在冰层中发出的灵质辐射。在霜脊诸神献祭后，他们残留在山脉中的那些霜纹碎片——那些曾经刻在瓦尔德林眼眶周围的、包含着古老智慧的符文——逐渐深埋到冰层中。冰本身不保留灵质信息。但那些曾经被符文的边缘割过的冰晶——它们记得一道边缘的触感。就是这种微弱到即使最敏感的帷慕学者也未必能察觉的"冰晶对一道古老划痕的物理记忆"——被某些植物学会了去捕捉、去吸收、去活。}


\textbf{这些植物于是变成了一种极端缓慢的、以世纪的尺度来计算的过滤器。它们从冰层中吸取到的神力残留——经过五、十、十五代的自然选择——在植物体内形成了各种不同配比的灵质化合物。有些化合物能被用来强化灵焰（埃洛伊丝的祖母活了一百四十岁——不是靠任何魔法，而是靠每天一剂特制的霜苔茶来轻微增强灵焰的天然屏障）。有些化合物则引发了完全相反的效果——它们会对蚀痕产生一种排异性反应。不是击杀蚀者——是让蚀痕"觉得这个宿主不好吃"而主动脱落。当然，"主动"这个词是拟人化的——蚀痕没有意志，它只是一个寄生性灵质结构。使蚀痕从一个特定宿主的灵焰上脱落——需要让那个宿主的灵焰瞬间产生一个极高频的灵质脉冲。而霜苔的一种特殊变种——"焚霜苔"——在晒干、研磨、掺入由骨灰平原上收集的被遗忘之名的粉末后——能精准地触发这种脉冲。埃洛伊丝在大量实验中发现的不是一蹴而就的"毒药"——而是这种脉冲的频率匹配。花了几百次尝试，她才从一个非常特定的配比（两份焚霜苔、一份白苔藓、三分之一份骨灰——不多不少）中得到了可靠的蚀灭反应。}


\textbf{但蚀灭药剂不是没有代价的。脉冲在摧毁蚀痕的同时也在摧残宿主的灵焰——那颗原本纯净、但已经被蚀痕侵蚀了七成以上的灵焰核心在脉冲中也会解体。灵焰中的原主残存意识——如果还有的话——在蚀灭反应中化为灰烬。所以你每一次使用蚀灭药剂，都在杀死两个存在：蚀者（已经不能被称为"人"了），以及还困在蚀者灵焰中的那个早已被淹没了的原主意识。埃洛伊丝临死前对自己杀了太多原主意识深感自责——她在笔记的最后一页写了一句被后来的每一代草药学者默诵的话："如果你的目标是帮助——你需要接受一个事实：有时你能给予的帮助只是送别。"}


\textbf{愈灵汤剂则更复杂。它不涉及蚀灭脉冲——它是通过极微量的灵植汤液（霜苔温和变种、混入圣橡落叶的浸液、再加一滴由灵织者采集的"命焰露"——那种在黎明时分从灵质层面凝结在草叶上的无形露珠）来修复灵焰上的细微裂缝。原理不是"愈合"——是人造的灵质填充物——它把灵焰上裂开的缝隙用植物的灵质纤维暂时性填充，让自然的愈合过程有时间发生。汤剂的有效时间是三天。如果三天内宿主自身的灵焰恢复力没能完成修复——汤剂中的纤维会被灵焰的天然免疫系统排出去——裂缝重新出现。所以愈灵汤剂不是最终解药——它只是"缓冲"。它给灵焰多三天的时间。有时候三天就足够了。有时候——当裂缝太深——三天只是将死亡推迟了三天。}


\section*{第三章：触碰之术——愈灵师的传承}


\textbf{愈灵师的能力是所有守幕者中最稀有的。它是天生的，无法通过训练获得。}


\textbf{在帷幕之域的灵质科学中，有一种叫做"灵焰触觉"的感知层次。普通人能感知到最粗糙的灵焰触觉——母亲安抚啼哭的婴儿、恋爱中的恋人对视时感受到的心跳加速、一对老夫妻在夕阳下并肩坐下时那团不必说出口的"我在这里"——这些都是最原始也是最普遍的灵焰触碰。你不需要任何训练就能感受到——你只需要是一个活着的人。但能将这些被动的、不自主的触碰转化为主动的、精确的、可以修复灵焰裂缝的诊疗——这种能力在人类中的发生概率大约是十万分之一。而且在每一个拥有这种天赋的人身上，它都会以不同的方式呈现——有的人能看到灵焰的裂缝（像墙上的一道裂纹），有的人能用手指触摸到裂缝的温度（热的是正在恶化的裂缝，冷的是已经稳定下来的老伤），有的人能在睡梦中自动进入伤者的灵焰中进行修复（而醒来后完全记不清修复的过程）。}


\textbf{奥尔德温——第一位愈灵师——的天赋是"灵焰联觉"：他能在灵焰触觉层次中同时感知到五感的信息。裂缝在灵焰中的形状他"看到"为一块撕裂的布；裂缝的温度他"摸到"为凉得烫手（因为灵焰的冷和热是反过来的——损伤的地方吸热，所以摸起来其实是热的，但体感是冷的）；裂缝的"声音"他"听到"为一道低沉的、几乎处于人耳能感知的极限低频的嗡鸣。所有这些信息同时涌入奥尔德温的灵焰——他根据这些信息的组合模式来判断裂缝的类型、严重度、形成时间、可能的扩散方向。然后他用自己的灵焰"触模"目标灵焰——不能用力过猛（否则裂缝被撕大），不能太快（否则目标灵焰的免疫反应会把他的灵焰当作外来入侵而攻击他），不能太久（超过五息的连续触模会使他的灵焰被感染）。在所有这些限制条件下，他必须在一息之内精准定位裂缝的位置并用自己的灵焰去"填充"。}


\textbf{这比做一场心脏手术更难——因为没有手术刀、没有X光机、没有任何可以放大和稳定的视觉辅助。你唯一能看到的是一团模糊的淡金色的火——而火中那些黑暗的裂缝在以不可预测的节奏和方向扩展或缩小。奥尔德温在四十年间做了数百次手术——成功率——据后来的统计——大约是三成四。不是他不合格——在当时，在没有任何前辈经验可参考、没有任何理论指导、没有任何训练体系的情况下——他全凭直觉维持三成四的成功率，是一个奇迹。今天的新一代愈灵师——在奈莉娅的灵焰影像技术（一种用灵质摄像机将灵焰的瞬时状态固定在一块水晶板上的工具）的辅助下——成功率约为七成。他们每年都会在奥尔德温的雕像前放一束刚从黑森林摘的枯枝。花已经谢了。但这不重要。重要的不是花还新鲜不新鲜。重要的是每一年都有人来到那尊雕像前。放下新的枝。这个仪式叫做——"你不会被忘记。"}


\section*{第四章：光盾之术——帷幕守卫的传承}


\textbf{马尔库斯的第一面灵焰护盾不是练出来的。它是被裂隙逼出来的。}


\textbf{他十八岁那年——那时他还是一个普通的守夜少年，在哨塔上打瞌睡——帷幕突然出现了一次五百年一遇的大规模波动。一道力场——后来被奈莉娅命名为"帷幕浪潮"——扫过了整个聚落。浪潮本身不含蚀痕——它只是帷幕在呼吸中的一次"深呼出"。但对灵焰来说——一次帷幕深呼出所携带的、来自神织物质的巨大灵质波动——像是有人在你耳边突然敲了一口大钟。近距离承受的人——如果没有任何防护——灵焰会被震出躯体。}


\textbf{马尔库斯当时正在哨塔上。他被力场直接从塔顶掀到了地面的石台上。他的两个同伴——就站在他旁边——当场灵焰离体。死亡。而马尔库斯在坠地的瞬间——他自己的灵焰在他的求生本能触发下做了一个他之前从未尝试过的事——将灵焰向外延伸了一道极薄的屏障。不厚——大概比一层宣纸更薄。但这层屏障在他身体周围形成了一个灵质绝缘层：帷幕浪潮的频率被他屏障的微频率部分抵消了。力的方向被偏转了。他活了下来。他的两个同伴没有。}


\textbf{之后的四十多年里，马尔库斯每天都在磨练他的灵焰护盾技巧。他发现他能将护盾投射到他人身上——只要那人与他有过灵焰触模。灵焰护盾不是被动技能——它需要守卫者持续地将自己的灵焰集中在目标身上。这不是一件容易的事，在夜晚、在恐惧中、在不知道目标是谁、不知道威胁从哪个方向来的情况下，要维持一个稳定到足以挡下一击的护盾——是对一个守卫精神集中力的终极考验。}


\textbf{而每一次护盾被击穿——守卫自己承受一部分灵焰冲击。这种冲击不会直接侵蚀灵焰——但它能在灵焰中撕开微小的结构裂缝。这些裂缝不是蚀痕（裂缝本身不会扩散——但它会使裂缝处的灵焰变薄）。变薄的灵焰更容易被未来的冲击击穿。所以每一次成功的守护——虽然救了被守护者——都在让守卫者自己变得更加脆弱。这就是"守卫重伤"——"灵蚀重伤"——的本质：不是被蚀痕感染了，是灵焰被反复撕裂后形成了太多连愈灵师都无法填补的裂缝。}


\textbf{马尔库斯在自己重伤的那几次——他的灵焰裂缝的数量超过了愈灵师能填补的上限——他不是躺在哨塔里等恢复。他继续守护。他会选择一个特定的庇护所——那个里面住着一个他从来没有见过的人——因为那个人的灵焰频率（他通过灵焰感知能读到）在附近所有频率中是最微弱的。那个人活下来的概率——如果不被守护——是零。所以马尔库斯每个夜晚都把他的最后一段完整的灵焰——那一段还没有被撕裂过的——延伸向那个庇护所。一年后，那个人搬出了聚落。马尔库斯不知道他为什么离开。但他收到了一封由一位行路灵织者代笔的信——信纸是从他自己随身背的一卷牛皮纸上撕下的。信只有六个字：}


\textbf{"你守护的人是自由的。"}


\textbf{最后一个字——"的"——在牛皮纸上被墨迹轻度洇开。那不是墨水——那是写信的人在写下这个字时落了一滴泪。那一滴泪在几百年后——当这封信从一个收藏家传到另一个收藏家——已经干成了牛皮纸上一个接近透明的水渍。但如果你对着烛火看——你能在那个水渍中看到一个很小的、模糊的、长得像送信人那双粗糙的大嘴唇的印记。}


\section*{第五章：灵痕之道——追猎者的传承}


\textbf{塞拉斯在黑森林中独自长大的头十年——大约从他六岁到十六岁——他是一个不说话的孩子。不是生理上不能说话——是没有人可以对话。黑森林中有鸟，有鹿，有刺猬，有蜘蛛，有一整个季节在林中蜕皮的蝉。但这没有一个是人。塞拉斯学会了一个人过日子的全套技能——用石头和树枝磨制箭头、在雨中挖坑储水、在枯死的圣橡树心刨出一个足够深的凹槽来躲避冬夜——但他的心智和社会交往的发展水平在他十六岁进入暮色聚落时还不如一个六岁的聚落儿童。}


\textbf{但有一项技能是他在森林中磨出来的——他在追踪动物（和后来追踪蚀者）时，逐渐发现了大自然中的灵痕。不是眼睛看的——他的眼睛和普通人一样只能看到普通的光线。他是在另一种层次上——一种他无法用语言描述的、比他任何其他感官都更"深"的知觉——看到了灵焰留下的印记。一个活人在某个地方走过——在他的每一只脚落地时，灵焰的最外缘会与草叶表面微微触碰。这层触碰在灵质层面上短暂地打开了一个很薄很薄的"印记"——不是物理上的脚印，是灵质层面上的"温度"变化。塞拉斯感觉那个印记像一块曾经被坐过但随后又变冷了的光滑石头——而那个石的凉——比其他任何地方的凉都要微微——一个他不能说出单位的微分——轻一些。这种感知太细腻了，细腻到他自己也无法确定到底是真实的还是他想象出来的。但当他每次沿着他"感觉"到的方向走——他最终总是找到了他想追的猎物。}


\textbf{进入聚落后，塞拉斯被一个老帷幕学者发现了。那个老学者在观察他追踪一只丢失的羊的过程中，发现塞拉斯追踪的不是羊的足迹——而是羊主人走过的路线上留下的、仍在微光中的灵痕。这种追踪能力——老学者意识到——可以成为一种碾压蚀者隐蔽能力的工具。因为蚀者在灵焰感知中是几乎隐身的——普通灵织者只能感知到"有些东西在外面"而不能定位。灵痕追猎者不需要感知蚀者本身——他只需要感知蚀者前几天留下的灵痕。蚀者的灵痕和普通人的灵痕在颜色、温度、厚度上完全不同——黑色的、冷的、比普通灵痕更"尖锐"（塞拉斯的词——他后来学会了足够的词汇，是他用每天一个新词的速度花了二十几年一点一滴地学起来的）。}


\textbf{塞拉斯的灵焰弩的制作过程在他死后很久被他的一个弟子公开。弩身材料是黑森林中一棵估计有两千岁的枯死圣橡的树干心材——那树干心材的灵质"记忆"中储存了这棵橡树生前从土中到土壤的、年复一年不变的、每一次季节更替中它通过根语与其他圣橡交换的安慰。塞拉斯说弩身不需要任何神力——圣橡自己的灵质残响就足够强。"它不是一根死木头。如果我在发射弩箭时闭上眼——我能听到它还在微弱的根语中和其他的树在说话。那些树已经不在了。但它还在叫它们的名字。"}


\textbf{发射不是箭——是将使用者自己灵焰中的一小片——极小的——可能只有整团灵焰的万分之一的碎片——从弩弦上弹出去。击中蚀者的灵痕后——击中点会像星火燎原一样沿着蚀者的整个灵痕网络迅速蔓延。蚀者在这个瞬间失去他的"方向地图"。他变得看不见任何方向的灵痕——像一个航海者突然失去了北极星的定位。他会走失——走到一个完全不认识的地方——然后发现自己无法回去了。晨曦将至。第一道霞光触到他时——那残留着拉赫卡里昂日舟远古神力的光——会开始烧毁他的反灵焰。烧得很慢。他会一直在那道光中站着——不能动——因为他的灵痕网络已经被击碎了，他不知道要逃往哪个方向。他只能在晨曦中等着——等着光的烈焰将他的一切归属空白。}


\textbf{塞拉斯的噬灭短刃是他的秘密武器。他从骨灰平原的一个被掩埋的古老锻造坊中找到了一把材质无法辨识的深黑色金属匕首。经过奈莉娅的分析——那不是金属，是冥界天秤崩塌后的碎片——安努赫卡审判亡灵时所用的天秤基座的一部分。秤座的石材本身不是"魔法"——但它在一个特殊的条件下（被安努赫卡在每一次审判后触摸过的两万年间）在微小量级上继承了天秤的原理：它和灵焰之间的重量关系是反过来的。正常的物体对灵焰的作用为零——灵焰和物理实体之间不会相互作用（这也是为什么蚀者可以穿越石墙）。但这把短刃——因为安努赫卡触摸过的每一道天秤都曾经在灵焰和纯白羽毛之间做过重量比较——对灵焰有了一种微弱的"物理吸引力"。当你用这把短刃接触一个蚀者——不是刺他的肉体——而是刺入他灵焰的边界那一刻，短刃会开始把蚀者灵焰中的蚀痕拉向自己——像一个磁铁吸铁粉。蚀痕被从蚀者灵焰中抽出、附着在短刃上。短刃在这个过程之后会暂时无法使用——它的表面被蚀痕覆盖了——需要用纯净的灵焰去"灼烧"（愈灵师的标准清洁流程）才能复原。而蚀者在那之后不是"死了"——是被"洗了"。灵焰中所有的蚀痕都消失了。但在这同时——如果蚀者原来还有一丁点自己的残余意识——那个残余意识也要直面一个事实：他在失去蚀痕的瞬间——突然重新"感觉到"了。他感觉到了他的手在哪里。他的皮肤上有风在吹。他的脚底接触着地面。然后——他想起来他曾经是人。}


\textbf{有些蚀者在被噬灭短刃刺后几息之内直接死亡——因为他们自己的灵焰早已被蚀痕彻底替代——蚀痕被移除后，没有任何灵焰剩余。有些蚀者——极少——在被刺后完全恢复了纯净。但他们的记忆——那些在被蚀变期间所做的事——那些被噬灵的人的面孔——那些在夜晚中发出过的无声尖叫——这些东西都留在了他们纯净化后的灵焰中。大多数恢复者在恢复后的几天内自己了结了自己。不是因为惩罚——是因为无法与自己共处。}


\textbf{塞拉斯在他最长的狩猎之后——他用噬灭短刃净化了一名已经噬过十七人、包括他自己的妻子的蚀者——在回到聚落后，他独自在哨塔上待了三天。没有人敢上去问他任何事。三天后，他下来了。他走到他的弟子面前，把噬灭短刃交给了他。他说了一句话——语气平静但每个字以重音结尾："他哭了。在最后。"}


\textbf{塞拉斯死于五十三岁。死时脸上带着一种说不清的表情——不像悲伤，不像解脱，不像任何可以被命名的感受。他的弟子们后来在整理他的笔记时发现了一张纸——纸上只有四个字："够了。给下一个。"}


\section*{第六章：灵织之道——织幕者的技艺}


\textbf{灵织者是守幕者中最大的群体——占全体的大约六成。他们既不是最强大的、也不是最善于战斗的——但他们是数量最多的。而在这个以灵焰为核心兵器的世界里，数量本身——如果被正确组织——是一种惊人的力量。}


\textbf{灵织者的共同基础是三种能力。第一种是}**\textbf{灵焰感知}**\textbf{——感知周围人的灵焰波动。不是像帷幕学者那样精确到蚀痕边界——而是更粗但更广的感知：谁今晚出门了？有几个人的灵焰在凌晨时分突然剧烈波动（可能是被攻击了）？谁的灵焰颜色从金色变成了灰白（可能需要愈灵师关注）？灵焰感知是一种先天禀赋加后天训练的混合技能。禀赋好的人，感知范围可以覆盖整个聚落。禀赋一般的人——感知范围大约是一至两个屋子半径。但可以练。有些人在二十年持续的夜行和守夜中，可以将感知范围从最初的四十几步拓展到整个聚落的边缘。}


\textbf{第二种是}**\textbf{帷幕低语}**\textbf{——在帷幕裂隙附近捕捉信息碎片。帷幕本身是由诸神残留意念编织而成——那些残留意念中包含的信息是海量的、杂乱无章的、没有标签的。但偶尔——如果你在裂隙附近、在血月之夜、在灵焰感知全开的情况下——你能"听到"一道很细的、像远处有人在为你耳语的信息。低语的内容可能是很多年前的某个夜晚、某个与你无关的人的一次选择。也可能——在极少的情况下——是非常近的、也许是昨晚某个人的灵焰残响中留下的最后一句话。帷幕低语不可靠——因为低语中混合了诸神残留意念中的随机噪声——你需要大量的交叉比对才能从中提取准确的情报。但它是灵织者们最原始的、不可替代的情报来源之一。}


\textbf{第三种是}**\textbf{灵织术}**\textbf{——将大量灵焰感知数据和帷幕低语碎片系统性地整合、交叉比对、编织成可读的判断。灵织术的名字来源于第一位织幕灵织者——那个用织布来比喻灵焰流动的女纺织工。她将灵焰的消长规律比作经线（不变的基础线）和纬线（每次波动都要穿过的、不断前进的新线）的交织——织出来的布上能看到灵焰的热点、冷点、异常点和可预判趋势点。今天，灵织术已经发展为一种高度系统化的技术——灵织者们在一张大牛皮纸上用灵质墨水（掺入自己的一滴灵焰残液——这滴残液赋予了墨汁在灵质层面发光的特性）绘制精确的"灵焰地图"。这张地图每晚刷新一次——由当夜的灵织者团队根据最新的感知和低语数据更新。新地图画好后——旧的归档。奈莉娅的观测塔地下室里保存着从最初到现在每一夜的灵焰地图——厚厚的一大叠。任何后来者都可以翻阅这些地图，看到聚落灵焰在千百年间的变化趋势。有的人翻阅之后大哭一场——因为他们看到自己的灵焰频率和某个已故之人的灵焰频率在地图上的线条从一开始离得很远，到后来逐渐靠拢，到最后完全重叠——这是一对终生相依、终于在死亡中融合了的恋人。}


\textbf{八种子类型的灵织者各有自己独特的感知侧重。}


**\textbf{老兵灵织者}**\textbf{——那些年轻时是追猎者或守卫、年迈后转为灵织者的老者。他们的灵焰经过几十年的战斗已经不再旺盛——但他们补偿以极其敏锐的判断力。一个老兵灵织者能在三息之内通过"气息"判断一个夜行者是否是被蚀痕沾染的——不是用察灵，而是靠一种更原始的直觉。老兵会说："那个人——闻起来有霜烧的味道。"霜烧是永冬山脉一种罕见的自然现象：霜晶在冬季被强风刮到岩石上，碰撞出微小火花。火花本身不够点燃任何东西，但它的烟雾有一种特殊的呛人气息。老兵灵织者用这个名字来形容蚀痕在灵质层面的"气味"——因为它像霜烧的烟雾一样——不是你"闻到"的，是你从脖子后面一阵突如其来的麻痒感觉到的。}


**\textbf{行路灵织者}**\textbf{——那些走南闯北的人。他们习惯在陌生的、危险的、随时可能有噬灵者出没的地方夜行。一个行路灵织者一晚上能进行两次独立的灵焰感知行程——在聚落两端之间快速移动。他们的灵焰感知范围不如一些固定不动的守夜灵织者——但他们的移动性使他们能收集到不同区域的灵焰数据，交叉比对后往往能得出更全局的判断。代价是——他们在外面暴露更久，被蚀者或其他危险碰到的概率是其他灵织者的双倍。行路灵织者中有很多人腿上都有旧伤——不是蚀痕，是被荒原上的瘴气、泥沼、或夜间出没的野兽留下的。但他们仍然每晚出门——因为没有人能替代他们走两次路。}


**\textbf{学徒灵织者}**\textbf{——那些天资不足成为愈灵师但仍想学习灵焰修复的人。他们的治愈能力非常基础——延缓死亡，不能逆转死亡。他们随身携带一个小药包——里面有几簇从聚落后山采的药草（迷迭香、百里香、鼠尾草——都是在帷幕之地随处可长、被灵质稍微沁染过的普通植物），用一条用旧了但被愈灵师祝福过的白麻布包着。他们把这包药放在濒死者的胸口——不是物理意义上的放置——是灵质层面的：将药草中的微弱灵质残余沁入濒死者的灵焰裂缝中。这最多只能推迟灵焰消散一蚀月（有时只有几天）。但这几天——对某些人来说——够他把最想说的话说完。}


**\textbf{记述灵织者}**\textbf{——那些有文化的人。在帷幕之地，文字是珍贵的。大多数聚落居民是文盲。但总有少数人——在他们的家族中流传下了从诸神时代继承来的文字能力。记述灵织者能在聚落会议上发挥无可替代的作用：他们不仅能将其他灵织者提供的原始感知数据用通俗语言解释给所有人听——还能用故事把复杂的灵焰动态转化为直观的道德抉择。"昨晚有三个人的灵焰在东北方向波动。一个是黑的——蚀者——移向中央广场方向。另外两个是金的——纯净但微弱——从同一个方向移向远离广场的庇护所。推测——这是一次未遂的噬灵。有人保护了这两个金的人。"这是一个枯燥的灵焰数据报告。记述灵织者会把它重新讲为："昨天半夜，有一个老人带着两个孙子跑过了聚落。追他们的东西很黑。他们差点就被抓住了。但在某个我们没有记录的角落——有一个帷幕守卫或灵痕追猎者——把他们藏了起来。我们不认识这位藏他们的人。但他阻止了一次可能影响三个人的噬灵。我们欠他一顿饭。"——说完后聚落中没有一个人不知道该感谢谁、该警惕谁。}


**\textbf{守夜灵织者}**\textbf{——那些天生警觉的人。他们不需要刻意去感知——他们的灵焰自动捕捉周围灵焰波动。一个守夜灵织者能在半睡半醒中分辨出自己门外的脚步声是几个人——轻重如何——方向是去往哪里。守夜灵织者们也负责维持聚落的夜间秩序——他们轮流值班，两两一组，在不同的哨塔之间巡逻。如果有人在夜晚需要帮助（比如生病、临盆、或灵焰突然异常），守夜灵织者是第一个知道的——因为他们会在灵焰感知网中看到一团突然变得不稳定的金火。}


**\textbf{炊火灵织者}**\textbf{——那些在聚落中有最广泛社交网络的人。他们经营着聚落的几间共同的厨房和烘焙房。每个黄昏——在夜行开始之前——聚落中的人会在炊火灵织者的厨房外排成长队，领取一份由当地谷物和灵质沁染过的水烹制的简单晚餐。这不是奢侈品——是必需品。因为那碗饭中含有炊火灵织者在烹饪时注入的极微量灵质暖意（他们灵焰中的热度被转移到炉火中）。这层暖意能在夜晚为食用者提供非常微弱的额外灵焰防护——不是能挡下噬灵，但能让灵焰在遭受冲击时更有韧性。炊火灵织者也通过这个日常仪式建立了一个庞大的"食客网络"——他们知道所有人的灵焰频率（因为在煮饭时他们的灵焰与每个人的食物发生了间接接触），所以他们能在灵焰感知网中更精确地定位每个人的实时灵焰状态。这种精确度有时甚至超过了帷幕学者的察灵——因为学者的察灵需要刻意针对某个人，而炊火灵织者对所有人的"大致"状态都是同时感知的。}


**\textbf{锻甲灵织者}**\textbf{——那些能打造灵质防护器具的铁匠。他们锻造的不是普通的铁——是矿脉山脉深处一种含有灵质记忆的金属（那种金属在诸神时代曾是用来铸造霜脊圣座城门的材料——虽然现在的纯度和数量都不如从前，但依然保有对灵焰的部分排异性）。锻甲灵织者将这种金属锤打成巴掌大小的金属片——贴在庇护所的门内侧。当蚀者经过时——他的反灵焰会与金属片中的排异性产生微反应——金属片会轻微发热（约人体体温）——热到足以唤醒正在睡觉的屋主。这不阻挡蚀者入内——但给了屋主几息的时间做出反应：呼救、逃跑、或者激活自己身上的最后一道灵焰护盾。锻甲者自己每次锻造一件灵质门锁时，要在锻铁时将自己的灵焰沁入铁中——每一次锤击都在铁块的灵质结构中留下他灵焰的一个微小的、不能再被收回的印记。这就是为什么很多锻甲灵织者在五十岁后灵焰就已经薄如蝉翼——他们给了太多出去，留给自己太少。铁匠铺外的木架上钉着上百块用旧了的门锁——每一块门锁后面都附了一小片皮纸，写着门锁使用者的名字。有些名字已经被划掉了——不是因为他们还了锁，是因为他们已经不在了。但他们用过的锁——锻甲者不舍得熔掉——还在木架上等着。等一个也许永远不会再回来的人。}


**\textbf{织幕灵织者}**\textbf{——最罕见的灵织者子类型。其他灵织者感知灵焰——织幕灵织者能直接感知帷幕本身。他们的灵焰频率与诸神编织的帷幕之间存在一种罕见的共振关系——没有人能预测某个人是否会拥有这种共振，因为它与血脉无关、与训练无关、与灵焰强大程度无关。它只是——出现。一些平平无奇的人——在某个平平无奇的夜晚——突然发现自己能听到帷幕中的所有低语同时被同步为一个可理解的、完整的句子。他们把这种状态称为"幕合"——在那一瞬间他们不是在"听"帷幕——他们就在帷幕中。他们知道——在那一瞬间——这张由神织纤维构成的巨大网络中每一根纤维的振动。当"幕合"结束后（通常持续不超过七息），织幕灵织者会扑到纸上，疯狂地写下他们在幕合中听到的主要信息。写完后他们需要静卧一整天——体力和灵焰消耗都太大了。他们写下的内容——被后来的灵织者团队分析、交叉比对、提炼出可用的情报。其中最有价值的情报——被称为"幕谕"——是帷幕本身作为由诸神残留意念构成的集体意识发出的"意图信号"。幕谕不能预测未来——但它能告诉你帷幕当前正在"关注"什么。有时它会关注某个特定的人的灵焰——可能因为那个人的灵焰正在被蚀痕侵蚀——帷幕还没有决定要把这个灵焰拉到哪一边。有时它会关注聚落中的某一次选择——可能因为那次选择的结果会在帷幕的呼吸节奏上留下一个永久性的波动。}


\section*{第七章：灰域之道——冥僧人的悖论}


\textbf{冥僧人——蚀者中最独特的存在——代表了守幕者认知中最深的灰色地带。}


\textbf{一个冥僧人不是生来就是的。他不是被蚀变的——他主动选择了这条路上最危险、也是最有可能保留一丝自我的分岔口的。当一个普通人的灵焰有大约六成被蚀痕覆盖——剩余的约四成纯净灵焰还在苦苦支撑——这个人面临三种选择。第一：放弃抵抗，让蚀痕完成不可逆的扩散——几天内变成蚀者，失去所有意识和自我，只剩下以噬灵为本能驱动的反灵焰存在。第二：找愈灵师进行治疗——如果剩余的纯净灵焰面积还足够大——愈灵师可以尝试蚀痕净化仪式——成功率低（约三成），而且即使成功（蚀痕被净化），灵焰中也会留下永远无法愈合的结构性疤痕——此后的日子里，这个人的灵焰会比之前更脆弱——任何一个微小的打击都可能再次撕裂旧疤痕，甚至堕入不可逆转的更深一层蚀变。第三——只有极少数人在走投无路时才能做出的决定——主动将自己灵焰中残留的、尚未被蚀痕覆盖的那一小片纯净灵焰——献祭给帷幕之外的某个存在。不是被吞噬——是主动献祭。目的是达成一个交易：用这一小片剩下的纯净灵焰作为"借款押金"，换取从帷幕之外借来的三种能力——裂隙引导（在特定地点开启小型裂隙来感知混沌中的活动）、堕化之力（将自己的蚀痕注入另一个人的灵焰中、引发对方也走向蚀变）、以及灵焰遮蔽（将自身灵焰中的蚀痕暂时隐藏、使自己看起来仍像纯净者）。}


\textbf{这三种能力中的每一种在使用时都有"利息"——即冥僧人每次使用它们时，都需要支付自己灵焰中那最后剩下的纯净部分中的一点点。每一次堕化他人——冥僧人自己离完全蚀变更近一步（因为他在把自己的蚀痕注入他人时也同时在扩展自身蚀痕覆盖的范围）。每一次噬灵——他偷取来的灵焰不能完全弥补他的借贷（借贷的利率是高于他一次噬灵的收益率的——奈莉娅计算过——大约是每噬五次灵才能抵偿一次债务的利息）。所以他是一个不断消耗自己的身体、灵焰、和残余人性去维持自己暂时不蚀变的悲剧存在。}


\textbf{但最悲剧的还不是利息，而是那个"交易"本身的根本缺陷：你向一个不可名状的存在借了东西——但这个存在从来没有向你解释过"还清"的条件。它给了你三种能力——但它没有说明你究竟欠了多少。而你不尝试你的每一次行动都在"还"吗？还了多少？你还欠多少？这些——一概不知。所以你只能继续活在——不是希望——而是拖延——中。把完全蚀变的日期从今晚拖到明天。从明天拖到下个月。从下个月——如果足够幸运——拖到某个可能你永远也抵达不了的远方。}


\textbf{这也是为什么守幕者对于冥僧人的态度是分裂的。有的人认为冥僧人只是蚀者中"延迟决定放弃自己"的懦夫——他们应该被归入蚀者阵营一视同仁。有的人认为冥僧人是被自己也无法控制的外部条件推上这条路的——在极端情况下仍有残存人性——在某些特定条件下他们仍然可以被"挽回"（但挽回的方法至今没有任何愈灵师能找到——因为冥僧人的灵焰中蚀痕和纯净灵焰的边界被那个不可名状的力量人为强化了——净化不了）。}


\textbf{记录在案的、自愿终结的冥僧人——就是将借来的所有能力"归还"给帷幕之外的存在、然后用自己最后一丁点残存的纯净灵焰去将自己的整个蚀变灵焰引爆——化为纯白的骨灰——只有约十余人。余烬是第六个。}


\bigskip\hrule\bigskip


\chapter{卷六：十五魂印}


\section*{序}


\textbf{在帷幕之域的每一个角落——从北境的永冬山脉的冰裂隙，到南端骨灰平原的灰白沙土，从东边雾隐禅院的石径，到西边黑森林深处的渡鸦冢——游荡着一些特殊的人。}


\textbf{他们不是神。他们不是蚀者。他们只是凡人——但他们每个人身上都携带了某种独特的、来自已死诸神或太古时代的印记。这些印记不赋予他们神力——诸神已死，神力已逝。它们赋予他们的是某种"天赋偏向"：一种在特定领域的超凡感知力，一种由特殊经历塑造的独特视角，或是一种在灵质层面上不同于常人的特殊灵焰结构。}


\textbf{他们被称为——}**\textbf{魂印者}**\textbf{。}


\textbf{每一个魂印者的印记可以追溯到一个已知的起源——一个诸神时代的遗迹、一个被遗忘的古老种族的血脉、或一次在裂隙中的濒死经历所引发的不可逆灵焰变异。印记本身没有名字——它是灵焰中的一个小小的、微微发光的、有时金色的、有时银色的、有时只在一瞬间出现的符号。不同的魂印者看到自己的印记的方式各不相同——有人在镜中看到自己灵焰中的符号（如果他们也是灵视者），有人在梦境深处被一个不完整的声音反复叫醒，有人在一生中唯一的濒死经历中——在意识还在身体和不在了身体之间那么一息——看到了印记的全部。}


\textbf{十五位——这是迄今为止被记录在帷幕之地所有已知魂印者档案中的独特印记载体。}


\textbf{以下为他们的传记。}


\textbf{（此处列出十二位已记录的魂印者——另有三位魂印尚在研究中，暂不公开。完整传记篇幅较长——每一位魂印者的个人史、灵焰结构、独特能力都在官方档案中有完整记录。游戏中选择表层身份时可查看各自的完整故事。）}


\section*{魂印者·霜脊谱系}


\subsection*{西格德——最后的霜脊猎手}


\textbf{永冬山脉最北的一道冰脊上，有一尊被冻结在奔跑姿势中的巨大冰雕。}


\textbf{冰雕的高度——按人类的尺寸——大约是正常的二十倍。冰雕的形状——是一个人形生物（两条腿，两条手臂，一颗头），但不是人。它的后腿比前腿长——适合于在极深的雪中跳跃。它的指端不是指甲——是利爪——每根爪长约一个成年人的手臂。它的嘴——即使被冰封着——仍然张开到一个不可思议的角度，露出了上下两排共六十多颗被磨成尖锥形的后牙。}


\textbf{这是一个}**\textbf{霜巨灵}**\textbf{——在诸神时代之前就存在于此地的古老种族的最后幸存者之一。诸神黄昏前，霜巨灵的族群有数百人。他们在冰原上过着与世隔绝的生活，从不同任何圣座发生交集。诸神黄昏中，裂隙在冰原深处出现——霜巨灵们没有神的能力来编织屏障来保护自己的领地。他们选择了一个更原始但同样有效的策略——将所有族人——从最老的到刚出生的——集中在一起，由他们的族长释放了一道足以冰封整个族群的远古寒意。所有的人在同一瞬间被冻结为冰雕。他们不能动，不能想，不能感知时间——但他们是"还在"的。裂隙无法侵蚀冰雕——因为冰雕中没有灵焰波动——没有可以被混沌存在感知到的存在的信号。冰雕——在裂隙看来——就是死物。而裂隙——奈莉娅说过——对死物不感兴趣。它只关心存在。死物不存在。}


\textbf{西格德没有见过任何活着的霜巨灵。他出生时——那个霜巨灵族群的最后冰川已经静静地在冰原上站了大约一万年。但他的爷爷曾在一个极夜里看到过奇异的景象——月光以近似直角的某个角度照射在冰雕上，冰层变得透明了一层——他能看到冰雕的胸膛中有一团微小的、银白的、极缓慢地沉浮着的光——那是霜巨灵的灵焰，仍然还活着，仍然还在冰层下——等着。等什么？谁也不知道。也许等有一天裂隙被封上——冰雕就会化——霜巨灵会继续他们一万年前中断的狩猎。也许那一天永远不会来。但知道他们还在这里——还在某个层面上活着——就够了。}


\textbf{西格德从他的家族继承了猎枪——不是灵焰弩（那是灵痕追猎者的专利），而是一把由霜脊山脉特有的永冻铁锻造的老式燧发枪。它在西格德家族中传承了七代。枪托上刻着七代猎人的名字——每一个名字都是用一种酸性液体（从冰原上的一种地衣中提取的）蚀刻在铁上，所以铁锈从那道蚀痕中向外渗，形成了一圈圈红褐色的锈环。最新的名字——西格德的——还没有生锈。它被刻上去才两年。刻它的人是西格德本人——在他的上一代（他的父亲）死后的第二天。他继承这把枪的同时也继承了一个问题——他的父亲在最后一次狩猎中追逐了一头在冰原上已搜寻了三年的巨兽——一头据说灵焰能遮掩自身不被任何灵视者发现的灵质冰原狼。他父亲追了它三年。终于追上。然后——枪响了。然后——没有人知道然后——因为西格德是在三天后在冰原上被巡山人找到的。他那时四岁半。他坐在他父亲的尸体旁边。他父亲的猎枪横在他的腿上。有一个连巡山人也未能理解的事实——枪管内有一股残存的烟——但烟没有从管口冒出来。它在枪管里一动不动地凝着——像一团被暂停在时间中的微云。那团烟——至今还在。}


\textbf{西格德从未解释过那个冬季到底发生了什么。他只在他年迈后——有一次在聚落的烹火旁——对一个守夜灵织者说了唯一一句话："我父亲没有开枪。是狼。"}


\textbf{没有人知道这是什么意思。猎枪只有他父亲在使用——如果不是他父亲开的枪，是谁？是狼？一头狼怎么会扣动燧发枪的扳机？那个问题至今没有答案。但西格德——在每一夜出门夜行时——都用一种极慢的、与他那巨大的身形完全不符的步速在冰原上走。他每走几步就停下来。他跪下，用手指按住雪面。他不说话。他只是——在听。听冰层下面一万三千尺深的地方——那个霜巨灵族长在他的冰封中每一次缓慢的心跳。那一跳——自当初冰封时——一分钟一次。心跳着的霜脊巨人——在等着有一天裂隙被封上。}


\subsection*{芙蕾维娅——华纳遗孤}


\textbf{麦野圣地——比霜脊更古老、比裂霆更柔和、比翠庭更安静的圣座——在诸神黄昏的早期纷争中被从地图上抹去。不是占领——是抹除。一股来自裂隙的力量将整座麦野圣地从时空中直接撕出来——像一个书仆把书中不想要的某一页撕下来团了用火烧。}


\textbf{但火在一瞬间之前——一个四岁的女孩被她的父亲从地上捡了起来、推进了一道随机传送裂隙。她掉在霜脊山脉的南麓——蜷缩在一棵枯死的世界树祖枝下——嘴里含着一片金色的麦叶。她哭不出来。不是不悲伤——是悲伤太大了——超过了一个四岁孩子的身体和灵焰能承载的上限。她灵焰中的某一部分——那个本该承载"失去故乡"这种感觉的那一小片灵焰——在那一刻因为无法处理如此大的丧失——自封了。那片灵焰把自己封闭为一个极其微小的金色球体——悬浮在她灵焰的中心。它还在——但它不能被触碰。触碰它——触碰那块封存了四岁小女孩失去家园的记忆——就会把二十二年前那场丧失重新撕开。}


\textbf{芙蕾维娅被瓦尔德林收养后，在霜脊圣座中度过了她漫长的成长岁月。她不是霜脊的神——她来自另一个圣座、另一种神族——但瓦尔德林从未区别对待她。她学会了霜脊的语言，学会了华纳族失去故土之前留下的一箱笔记中的麦野古语（她自学的——那些笔记用的是麦野圣地失传的灵质墨水——必须将灵焰温度调配到一个非常精确的水准才能显影），学会了如何在两种身份的缝隙中找到一个站得住脚的支点。她至今仍是霜脊遗民中血脉最远——但也最不能简化为任何单一血统的人。}


\textbf{她的灵焰是双色的——一半霜脊淡青，一半麦野金黄。在每年的夏至——当霜脊山脉短暂地进入两个月的融雪期——她会独自走下山——走到那片她从来没见过但在梦中不断出现的地方——一块山脚下的小草甸——这里的地形恰好和她从未亲身看过的麦野圣地有某种不可名状的相似——躺下去，闭上眼睛，试图以记忆中最原始的碎片重建那失去的一切。她一次也没有成功过。但她每年都试——这本身已是她与故土之间仅存的联系。不是地点——地点在物理层面上不存在了。不是人——所有的人都被抹去了。而是"尝试"本身——这个每年重复的动作——就是麦野圣地所有剩余存在的全部。当她躺在草甸上闭上眼时——她灵焰中的那片封存的金色球体会微微暖一点。不会解封。但会暖。}


\chapter{卷七：千面之影}


\section*{第一章：奈莉娅的最后观测}


\textbf{奈莉娅在她的观测塔上站在了第三十一年。}


\textbf{三十一年——她用一只眼睛（另一只在一次近距离裂隙观测中永久失明，瞳孔变成了固定在一半张开状态的、无法缩放的、但也因此获得了一种极其罕见的"裂隙透视"能力的替换品）每天观测帷幕上的每一条波纹、每一缕微光、每一丝颜色变化。}


\textbf{她的最重大发现——在她观测生涯的后期——是她在帷幕的呼吸节奏中筛选出了一个周期性出现的、规则的、不属于任何自然帷幕波动的低频信号。它很弱。弱到只有奈莉娅的敏感度才能从背景噪音中分离出来。它的频率——经过多年的记录和验证——与帷幕之外的某个具体存在的活动节奏一致。}


\textbf{那个存在——奈莉娅在《帷幕观测录》第七卷中第一次给了它名字——}**\textbf{千面之影}**\textbf{。}


\textbf{这不是它的真名——它的真名如果被任何人类获悉，会在听到它的人的灵焰中自动烧出一个无法修复的洞。这是奈莉娅从它的行为模式中提取的特征。它有一个独特的策略——不同于其他混沌存在。它不像其他存在那样在帷幕表面反复猛撞——这没有用——帷幕太厚了，诸神编织的纤维数亿万计的层叠。它也不试图在极短时间内积累强大力量冲破帷幕。它的策略——是渗透。}


\textbf{千面之影将自己的存在分割成了碎片——无数个微小到无法被帷幕整体探测的碎片。每一片碎片穿过帷幕上的微观裂隙——那些裂隙小到连赫尔米昂的金瞳都无法察觉——进入帷幕之域内部。每一片附着在一个蚀者的蚀痕中——不是控制蚀者（碎片太小了，没有完整意识的千面之影不可能控制任何人），而是"共享感知"。就像一个人把无数只微型的耳朵放在无数个不同的位置——千面之影的碎片在蚀者的蚀痕中被动地接收着来自蚀者周围所有灵焰活动的信息。恐惧、贪婪、爱、恨、牺牲、背叛——每一种人类情感产生的灵焰波动——都会被碎片采集、编码、通过碎片与千面之影本体之间的某种超越空间和时间的量子纠缠关系——瞬间上传到千面之影的"感知网络"中。}


\textbf{千面之影于是在两万年来通过这种方式感知到了人类的每一种情感。它品尝了爱——不是一场或一百场——是几十万场。它品尝了背叛——每一次被朋友放逐出聚落的尖叫、每一个在最绝望时出卖了队友的颤抖。它品尝了牺牲——那些在夜晚将自己的灵焰庇护延伸到一个自己根本不认识的人身上的匿名帷幕守卫——他们的每一次舍身——每一次——都在千面之影的网络中留下了一个它无法理解但很好奇的数据点。那是什么？为什么有人会将自己的灵焰撕碎来换取一个陌生人的存活？为什么有人类选择的行为模式在逻辑上完全逆反了他自身的生存本能？}


\textbf{千面之影不理解牺牲。不理解——所以它更加渴望去体验牺牲。但它是混沌的一部分——它无法进行牺牲。牺牲是一种只属于"存在"的行为——你首先要存在，你才能牺牲。千面之影还没有存在——它只是碎片式的感知网络。为了体验牺牲——它需要先存在。它需要——一个人。一个它在帷幕之内的"容器"。一个它将足够多的碎片注入足够长的时间、使碎片在那个人身上累积到足以产生一个微量的、跨维度的量子纠缠场——在那个人的灵焰中创造一个可以由千面之影的"分身"远程操控的小型存在接口。这个人——在奈莉娅的分类中——被称为"隙间者"——一个介于人与混沌之间的半人。隙间者不会意识到自己已被种下了碎片——但他们的选择模式发生着微妙的、不易被察觉的变化。他们比之前更容易做出不道德的决定。不是被强迫——是被提供了"一种更容易的选择"。在道德困境中——如果你听到了一个低语告诉你"走更容易的那条路——没有人会知道的"——你也许不会听。但如果你连续听了一年、每天、每个夜晚——在你最颓丧的时候、最绝望的时候、最疲倦的时候——那个低语都在。你会不听吗？}


\textbf{奈莉娅在观测塔的最顶层——在她被风吹乱的木桌上——画下了一个巨大的、由细小圆点连成的网络图。每一个圆点代表着一个她通过帷幕观测到的、带有可能的隙间者特征的灵焰。她还在继续添加新的圆点。她还没有完成这张图——因为她还没有找到那个中心——那个碎片累积最多的、最接近被千面之影选为容器的隙间者。但她知道——那个人在这张图的某处。}


\begin{quote}
\textbf{奈莉娅最后一篇观测笔记的最后一行：}

\end{quote}

\textbf{>}

\begin{quote}
\textbf{"我在找它。它也在找我。我们在同一片雾里。谁先看见对方——谁赢。"}

\end{quote}


\bigskip\hrule\bigskip


\chapter{尾声：织机之上}


\textbf{奈莉娅在观测塔的顶层放着一台小型织机。}


\textbf{不是她的——是第一代织幕灵织者（那个用织布来比喻灵焰流动的女织工）传给她的。那些曾经织入布匹中的、已死守幕者的灵焰残响——它们没有消散——它们留在了织机的木纤维中——像被冰封在琥珀中的微小的远古飞虫。}


\textbf{当奈莉娅把手放在织机上时，她能听到他们。不是用耳朵——是在灵焰触觉的最深的层次——听到那些已死之人的声音。他们从不发出完整的句子——只是碎片式的、在各种情绪层级上的微弱振动。有的碎片是暖的——那是某位已死守幕者在生前最后一夜的满意度（"我做到了"）。有的碎片是涩的——那是某位守幕者临死前的遗憾（"不应该那样做"）。有的碎片只是一个单字——"来"——没有任何上下文。奈莉娅不知道"来"是什么意思——来哪里——但她感觉那个单字是暖的。它是邀请，不是命令。它是一个人在另一个世界——也许是帷幕中、也许是无光之渊里——对她说："你准备好的时候——来。我们都在。"}


\textbf{奈莉娅在那台织机前独自坐了很多个夜晚。她把那些碎片听了无数遍。她学会了分辨谁是谁——哪一个碎片来自奥尔德温（他的碎片最暖——因为他在治愈了那么多人之后满足地死去），哪一个碎片来自马尔库斯（他的碎片中有一种冷冷的忠诚——即使在死后，他的残响仍然在守护着什么），哪一个碎片来自埃洛伊丝（她的碎片中有一片霜苔的苦味——那是她每天喝霜苔茶留下的残香）。}


\textbf{但她最想听到的那个碎片——来自余烬——来自那个选择将自己最后一点人性碎片归还给混沌、了结了自己的借贷、化为了纯白骨灰的冥僧人——她一直没有听到。她不知道余烬的声音在哪一层、怎样的频率中。也许他不在织机里——也许他没有灵焰残响留在任何地方——因为他已经在还清最后一笔债时把自己所有的碎片都归还了。他无债一身轻。他没有留下任何东西。}


\textbf{不——他留下了银链。那条刻着"我把借的都还了"字样的、在冥僧人化灰后掉在石头上叮响了一声的银链。那条链——后来被某个灵痕追猎者捡起戴在了自己的脖子上——然后传了好几手——现在挂在奈莉娅观测塔的朝北窗口上。西北风把它吹动了——它偶尔会轻敲窗框。敲窗的频率——奈莉娅数过——是每分钟五次。不是随机的。是一个人躺在床上——心跳缓慢、平稳、再不会突然加速——的心率。是一个已经不需要再担心任何债了的人——终于可以不分心去面对最想面对的那个瞬间——的节奏。}


\textbf{奈莉娅走到窗口。她把银链从窗框上取下，戴在自己的脖子上。链坠是一小块银片——只有指甲盖大——上面用很细很细的、明显是自学的刀法刻着一行字："我把借的都还了。"在字旁边——还有一些更小的、必须凑得非常近才能看到的——用极细的刻刀尖划出来的——也许不是字，而是某种冥僧人之间才懂的标记——一个小点带有五条线——代表什么？没有人知道。奈莉娅说——那是他留给后面的人的。如果还有后来的人选择走上这条路——如果还有人决定做相同的交易——这条链上已经有一个解答了。不是什么都告诉他——只是告诉他——"我曾在这条路上走到过终点。终点有光。"}


\textbf{奈莉娅站在窗口。面朝北——面朝着霜脊山脉的方向——面朝着她相信那些被冻结在冰层下面的霜巨灵仍然还在沉睡着、心跳一分钟一次、等待着裂隙被封上的那天。}


\textbf{她把银链攥在手里。她的灵焰——那只剩了三十年观测生涯、被裂隙灼伤过数次、已经不再旺盛的灵焰——在银链的灵质记忆旁边短暂地亮了一瞬。不是复活——只是打招呼。是在对链上那个早已不在的人说："我知道了。谢谢你。"}


\textbf{然后她回到织机旁。她开始整理她的最后一卷笔记。她知道自己这辈子不可能看到帷幕外的世界。但她可以把整理好的数据——三十一年来每一夜的观测记录、每一段被证实的推论、每一个被否定的假说——留给后面的人。也许在某个她还不会到达的时代里——会有一个女孩——也和她一样，坐在观测塔上——用那只唯一还完好的眼睛——看到一条她没有看到过的波纹——然后找到那个她这辈子也没能找到的答案。}


\textbf{"然后她会在笔记的空白处——写下她没有说完的句子。"}


\begin{quote}
\textbf{到这里。}

\end{quote}


\bigskip\hrule\bigskip


*\textbf{帷幕之地 · 灵焰纪元 —— 完结}*


*\textbf{但你的故事，才刚刚开始。}*


\bigskip\hrule\bigskip


\begin{quote}
*\textbf{"我把借的都还了。"}*

*\textbf{——余烬，冥僧人，最后一年，第十三天}*

\end{quote}


\end{document}
